\chapter{文本生成与转换}
\label{chap:nlp-text-generation}

前一章已经能够判断通知的类别、找到相关材料，并把重复报道组织到一起。现在需要交付的
不再是标签或材料列表，而是一段可以直接阅读的文字：把通知译给外国来宾，压缩成日历说明，
或改写成容易理解的参加指南。三种输出都可能写得流畅，却犯下不同的错误。译文不能把
“校外来宾”变成所有参加者；摘要是否必须保留报名截止条件，取决于它的用途；改写若删掉
“须”，就改变了通知的要求。

本章讨论怎样依据材料形成新文本，以及怎样控制其中的信息变化。条件生成提供共同的训练、
搜索和结果恢复接口；写作、编辑、摘要与翻译则分别规定内容从哪里来，什么必须保留，
什么可以改变。内容核验和评价贯穿这些任务，而不是生成之后附加的一项流畅度评分。
第~\ref{chap:transformer}章的编码与生成结构是模型基础，
第~\ref{chap:machine-learning-paradigms}章提供监督、预训练与迁移的学习背景。
本章只展开它们接入具体任务后的计算。

以下通知、资料、分数和轨迹均为教学构造，不是模型实测。共同通知于2025年4月15日
（北京时间）发布，活动编号为\code{ML-0418}：
\begin{quote}
机器学习读书会原定于4月17日举行，现改至4月18日（周五）14:00，地点为海棠楼201。
校内师生可直接参加，校外来宾须在4月16日18:00前提交报名信息。
\end{quote}
材料没有说明主讲人、费用、人数或时长，生成器不能把这些空白自动补成事实。

\section{条件生成的共同实现}
\label{sec:nlp-generation-interface}

形成文本不一定要逐词预测。模板可以填字段，摘要可以选择原句，纠错可以执行少量编辑。
这些路线都需要一份明确的输入、一种候选形成方式和一种验收办法。
图~\ref{fig:nlp-generation-interface}把三者放在同一接口下：训练参考用于估计参数，
运行时反馈用于定位候选中的错误，两类信息不能混同。

\begin{figure*}[htbp]
  \centering
  % Adapted from academic-drawing/templates/tikz_template.tex to the book palette.
\begin{tikzpicture}[
  every node/.style={font=\figurefont\small,align=center,line width=.6pt},
  block/.style={book node,text width=39mm,minimum height=12mm,
    font=\figurefont\small},
  flow/.style={book arrow,rounded corners=1.2mm},
  train/.style={flow,dashed,draw=FigureLoss},
  feedback/.style={flow,dash dot,draw=BookAmber}
]
  \node[block,fill=FigureInput!12,draw=FigureInput] (material)
    at (-51mm,0) {输入材料\\来源、对象与时点};
  \node[block,fill=FigureInput!12,draw=FigureInput] (task)
    at (0,0) {任务要求\\语言、篇幅与格式};
  \node[block,fill=FigureLoss!10,draw=FigureLoss] (reference)
    at (51mm,0) {训练参考\\目标文本与结束符};
  \node[block,fill=FigureModel!12,draw=FigureModel,text width=62mm] (candidate)
    at (0,-25mm) {形成候选\\模板、选择、编辑或序列生成};
  \node[block,fill=FigureModel!8,draw=FigureModel] (translation)
    at (-51mm,-52mm) {翻译检查\\原意、实体与否定};
  \node[block,fill=FigureModel!8,draw=FigureModel] (summary)
    at (0,-52mm) {摘要检查\\选择、支持与遗漏};
  \node[block,fill=FigureModel!8,draw=FigureModel] (editing)
    at (51mm,-52mm) {编辑检查\\允许修改与意义保持};
  \node[block,fill=FigureOutput!12,draw=FigureOutput,text width=62mm] (output)
    at (0,-78mm) {通过：恢复并交付文本\\失败：保留错误与证据};

  \draw[flow] (material.south) |- (candidate.west);
  \draw[flow] (task.south) -- (candidate.north);
  \draw[train] (reference.south) |- node[pos=.25,right] {参数监督} (candidate.east);
  \coordinate (split) at (0,-39mm);
  \draw[flow] (candidate.south) -- (split) -- (summary.north);
  \draw[flow] (split) -| (translation.north);
  \draw[flow] (split) -| (editing.north);
  \coordinate (join) at (0,-66mm);
  \draw[flow] (translation.south) |- (join);
  \draw[flow] (editing.south) |- (join);
  \draw[flow] (summary.south) -- (output.north);
  \draw[feedback] (output.east) -- (79mm,-78mm) -- (79mm,-29mm)
    -- node[below,font=\figurefont\footnotesize] {重新选择或修订}
    ([yshift=-4mm]candidate.east);
\end{tikzpicture}

  \caption{条件生成的共同接口。实线表示运行时的材料与候选流，虚线表示训练参考对参数的
  监督，点划线表示核验失败后的重新选择或修订。翻译检查原意，摘要检查内容选择与来源，
  编辑检查修改必要性；共享接口不表示共享同一训练数据或验收标准。达到修订预算仍未通过时，
  应返回失败及原因，不把最后一次候选视为合格输出。}
  \label{fig:nlp-generation-interface}
\end{figure*}
\FloatBarrier

\subsection{输入内容、要求与目标文本}
\label{sec:nlp-generation-training}

把输入条件记为$c=(x,u,b)$：$x$是待处理材料，$u$是任务要求，$b$是允许使用的背景。
例如$x$为上述通知，$u$可以是“译成英语，保留全部参加条件”，也可以是
“为校外来宾生成两句参加指南”。$b$可以给出年份、时区和经确认的名称译法，
不能偷偷包含测试参考答案。目标文本$y_{1:T}$采用固定分词器，其中$y_T=\mathrm{EOS}$
是结束符；起始符$\mathrm{BOS}$只用于启动生成，不是要交付的文字。

这一区分使同一材料的多个目标不再混乱。英文译文中的
“External guests must submit their registration information before 18:00 on April 16”
需要保留截止条件；只供日历显示的“读书会改至4月18日14:00，海棠楼201”
可以不承担完整参加指南的职责；面向来宾的易读改写则应写明
“校外来宾请在4月16日18:00前提交报名信息”。多参考文本表示同一要求下的合法表达差异，
不是把不同用途的答案放在一起，让模型任选一个。

为了学习由条件到文本的映射，采用\term{条件最大似然}：提高训练参考在给定条件下出现的
概率。链式法则把整段文本的概率分解成下一词元概率：
\begin{equation}
  p_{\bm{\theta}}(y_{1:T}\mid c)
  =\prod_{t=1}^{T}p_{\bm{\theta}}(y_t\mid c,y_{<t}).
  \label{eq:nlp-generation-factorization}
\end{equation}
这里$\bm{\theta}$是可学习参数，$y_{<t}$表示位置$t$之前的目标词元。
对独立训练样本的似然取对数，再改变符号，就得到最小化目标
\begin{equation}
  L(\bm{\theta})
  =-\sum_{i=1}^{n}\sum_{t=1}^{T_i}
  \log p_{\bm{\theta}}(y_{i,t}\mid c_i,y_{i,<t}).
  \label{eq:nlp-generation-nll}
\end{equation}
优化器据此调整参数，使正确词元相对同位置的其他候选获得更大概率。
先把每条样本的损失除以自身目标长度，再对样本平均，会让各句等权；
把全部有效目标词元的损失一起平均，则使长句贡献更多。
仅给同一个损失总和更换整体分母，不会改变样本间的相对权重，训练记录应写清归一化层级。

\term{教师强制}（teacher forcing）是在训练时把真实的前一段目标交给模型，而不是把
模型自己刚生成的错误继续传下去。设简化的目标词元为“周五／举行／EOS”，
解码输入便是“BOS／周五／举行”，预测标签仍为“周五／举行／EOS”。
若三个标签的条件概率为$0.6,0.5,0.8$，该样本的概率为$0.24$，
负对数似然为$-\log0.24\approx1.4271$，每个目标位置平均为$0.4757$。
结束符参与监督，模型才有机会学到何时停止。运行时只有已经生成的前缀可用，
训练与运行的前缀分布因此可能不同；增加搜索宽度不能消除这种差异。

编码器--解码器把$c$交给编码器，解码器读取移位目标，目标侧各有效位置计算损失。
仅解码器通常拼接“条件／分隔符／目标”，在同一因果序列内预测后继；
条件位置和补齐位置的损失掩码为零，目标及其EOS位置为一。
掩掉输入损失不等于掩掉输入注意：目标仍须能读取条件。
若一个任务采用不同掩码或同时训练输入语言模型损失，应明确说明改变了哪个目标。

训练材料可以来自人工平行文本、审订记录或明确标识的合成数据。提示示例与微调分别改变
运行时上下文和模型参数，其通用原理见第~\ref{sec:transfer-learning}节。
评价时允许看到的是$x,u,b$、已发布的词典和指定资料库，以及自身生成历史；
参考答案只交给评价器。选择提示、检索资料或阈值时读过的测试参考，不能再作为独立测试。

\subsection{候选文本的形成}
\label{sec:nlp-generation-candidates}

模型给出分布，应用还要从分布得到一个或多个具体候选。确定性搜索偏向高模型分数的序列，
随机采样保留多种表达的机会。两者都可以接在内容规划之后，也可以只负责模板中的一个槽位，
或局部编辑中的替换片段。长度、术语与格式约束见第~\ref{sec:nlp-generation-constraints}节；
候选重排则是在既有候选中重新比较，不会找回已经被搜索永久删掉的答案。

\subsubsection{贪心解码与束搜索}
\label{sec:nlp-generation-beam}

\term{贪心解码}（greedy decoding）每一步只选概率最高的后继。
\term{束搜索}（beam search）同时保留$b$个尚未结束的前缀，以降低早期选择失误的影响。
用累计对数概率$\ell(h)$给前缀$h$计分，扩展一个词元$v$时计算
\begin{equation}
  \ell(hv)=\ell(h)+\log p_{\bm{\theta}}(v\mid c,h),
  \qquad s(h)=\frac{\ell(h)}{|h|^\alpha},\quad \alpha\geq0.
  \label{eq:nlp-generation-beam-score}
\end{equation}
$|h|$不含BOS，完整候选的长度含EOS；空前缀只用于初始化，不代入归一化式。
本节用幂次长度归一化作为明确的教学约定，并非所有工具包的默认公式。
同一步的活动前缀等长，归一化不改变排序；不同长度的完成候选则可能换序。
由于对数概率非正，$\alpha>0$会减轻长序列的累计负分，但不能保证恰当长度。

\begin{algorithm}[带独立结束集合的束搜索]
  \label{alg:nlp-generation-beam}
  \begin{algorithmic}[1]
    \Require 条件$c$，束宽$b$，最大生成步数$m$，指数$\alpha$
    \Ensure 完成候选，或无完成候选的失败状态
    \State 活动集合$H\gets\{(\epsilon,0)\}$；结束集合$F\gets\varnothing$
    \For{$t=1,\ldots,m$}
      \State $C\gets\varnothing$
      \For{每个$(h,\ell)\in H$及词表中的每个$v$}
        \State $\ell'\gets\ell+\log p_{\bm{\theta}}(v\mid c,h)$
        \If{$v=\mathrm{EOS}$且满足结束条件}
          \State 把$(hv,\ell')$加入$F$，以后不再扩展
        \ElsIf{$v\ne\mathrm{EOS}$且前缀合法}
          \State 把$(hv,\ell')$加入$C$
        \EndIf
      \EndFor
      \State $H\gets C$中按$s$最高的至多$b$项；平局按预定词元序
      \State $F\gets F$中按$s$最高的至多$b$项
      \If{$H=\varnothing$}
        \State 结束循环
      \EndIf
    \EndFor
    \State 返回$F$中的最佳项；若$F=\varnothing$，报告预算内未完成
  \end{algorithmic}
\end{algorithm}

算法~\ref{alg:nlp-generation-beam}不把结束候选重新放回活动束，也不把达到长度上限的
未结束前缀伪装成完整文本。它最多执行$m$轮，每轮至多扩展$b|V|$个符号，
这里$|V|$是词表大小，模型前向计算的代价还需另计。
为便于核对，算法不采用提前终止。若$\alpha=0$，未来追加的对数概率不会增加累计分数，
当前最佳结束分数不低于所有活动分数时就可提前停止寻找最佳项。
若使用长度归一化，不能继续套用这个判断；有限$m$时，活动前缀的
$\ell(h)/m^\alpha$可作为此归一化下的乐观上界，但保证也仅相对于尚未剪掉的前缀。

\Needspace{12\baselineskip}
\begin{example}[三步生成中的一次早期遗漏]
  \label{ex:nlp-generation-beam}
  词表只有$A,B,\mathrm{EOS}$，$A,B$是抽象文字片段；最多生成三步，EOS占一步。
  下表给出所有会影响计算的条件分布，其他前缀统一采用$(0.10,0.10,0.80)$。
  \begin{center}
    \begin{tabular}{@{}lrrr@{}}
      \toprule
      已有前缀 & $A$ & $B$ & EOS\\
      \midrule
      空 & 0.55 & 0.40 & 0.05\\
      $A$ & 0.45 & 0.35 & 0.20\\
      $B$ & 0.10 & 0.10 & 0.80\\
      $AA$ & 0.35 & 0.25 & 0.40\\
      $AB$ & 0.10 & 0.10 & 0.80\\
      \bottomrule
    \end{tabular}
  \end{center}
  贪心依次选择$A,A,\mathrm{EOS}$，完整概率为
  $0.55\times0.45\times0.40=0.099$。
  取$b=2,\alpha=0$时，第一步活动束为$A:0.55,B:0.40$，
  空文本的EOS候选$0.05$另存。第二步从$A$得到
  $AA:0.2475,AB:0.1925,A\mathrm{EOS}:0.11$；
  从$B$得到$BA:0.04,BB:0.04,B\mathrm{EOS}:0.32$。
  活动束保留$AA,AB$，结束集合保留$B\mathrm{EOS},A\mathrm{EOS}$。
  第三步的结束候选分别为$AA\mathrm{EOS}:0.099$和
  $AB\mathrm{EOS}:0.154$；其余未结束扩展因预算耗尽不能交付。
  最佳完成文本为$B$，概率$0.32$。贪心的遗漏发生在第一步删掉$B$，
  不是第三步不会选择EOS。
\end{example}

该例中束搜索确实找到更高概率的完成序列，但$B$是否是好译文或忠实摘要，还没有任何
任务证据。模型偏好的短句可能漏条件，大束宽也可能更充分地暴露模型的这种偏好。
空文本虽然在该教学分布中合法，真实任务通常还有非空等硬条件；这些条件须进入结束检查，
不能在失败后临时改变规则来夸大搜索效果。

\subsubsection{top-k采样与核采样}
\label{sec:nlp-generation-sampling}

开放写作常常希望保留不同措辞或情节。直接从整个词表采样又可能频繁抽到低概率片段，
因此可以先调整分布，再限制候选集合。
\term{温度}（temperature）$\tau>0$通过
\begin{equation}
  q_\tau(v)=\frac{\exp(z_v/\tau)}{\sum_{w\in V}\exp(z_w/\tau)}
  =\frac{p(v)^{1/\tau}}{\sum_{w\in V}p(w)^{1/\tau}}
  \label{eq:nlp-generation-temperature}
\end{equation}
调整下一词元分布，其中$z_v$是未归一化分数。较小温度集中概率，较大温度使其分散。
$\tau=0$不能直接代入公式；唯一最大值下的零温极限才对应贪心选择。

\term{top-k采样}保留$q_\tau$最大的$k$项；
\term{核采样}（nucleus sampling，也称top-p采样）把词元按概率降序排列，
保留累计概率首次达到阈值$p_0$的最短前缀集合$V_{p_0}$。
两者随后都把保留概率除以集合内的总概率，再随机取一个词元。
核采样的集合大小随当前分布变化，不是固定保留$p_0|V|$个词
\scite{nlp-generation-holtzman2020}

沿用例~\ref{ex:nlp-generation-beam}第一步的$(0.55,0.40,0.05)$。
在$\tau=1$时，top-k的$k=2$保留$A,B$，归一化为$(11/19,8/19)$。
top-p的$p_0=0.50,0.90,0.98$分别保留$\{A\}$、$\{A,B\}$、
$\{A,B,\mathrm{EOS}\}$。若先设$\tau=0.5$，平方后归一化得到
$(0.65054,0.34409,0.00538)$；此时$p_0=0.60$只保留$A$，
而原温度下同一阈值会保留$A,B$。所以比较时必须固定温度与截断次序。

一次采样还必须说明何时停。这里每条候选最多三次取样，抽到EOS立即结束，
未抽到则标为截断；只在展示时去掉EOS，不在取样时删掉它。
固定输入分布、四条候选、Python的\code{random.Random(7)}，每种方法分别重置种子，
按降序累计概率用一次均匀数选词。top-k取$k=2$得到
“$AAA$（截断）、$AA$、$AA$、$AA$”；top-p取$p_0=0.98$得到
“$AAA$（截断）、$AB$、$AB$、$AA$”。两组$\tau=1$，最多各用12次取样，
本次均实际用满12次。若教学硬条件为“禁止相邻$A$”，前组没有合格完成项，
后组只有两条$AB$合格，且它们是重复候选。这个小样本只验证取样机制，
不能作为核采样优于top-k的实验结论。

记录唯一完成候选数、重复片段、截断率和条件满足率，比只展示最喜欢的一条更完整。
真实比较还需固定模型、分词器、提示、随机数实现与候选预算。
翻译也能使用采样，但多样性仅限于合法译法；温度不会把错误日期变成可接受的创作。

\subsection{结果解析、核验与修订}
\label{sec:nlp-generation-recovery}

模型词元不是最后的交付格式。恢复过程按分词器规则合并子词，处理空格和特殊标记，
保留段落边界；结构化输出则先解析字段名、类型和必填项，再检查值。
例如生成的日期字符串可解析为2025年4月18日，随后仍要核对它属于
\code{ML-0418}的活动时间，而不是报名截止时间。可解析与值正确是两项检查。

编辑任务还要知道哪里发生了改变。\term{Levenshtein对齐}使用插入、删除和替换恢复
原序列$x_{1:n}$到候选$y_{1:m}$的最少编辑。设$d_{i,j}$是两者前缀间的编辑距离，
以$d_{i,0}=i,d_{0,j}=j$初始化，递推
\begin{equation}
  d_{i,j}=\min\left\{
  \begin{aligned}
    &d_{i-1,j}+1,\\
    &d_{i,j-1}+1,\\
    &d_{i-1,j-1}+\mathbb{I}[x_i\ne y_j]
  \end{aligned}\right. .
  \label{eq:nlp-generation-edit-distance}
\end{equation}
三项分别删除$x_i$、插入$y_j$、替换或保留当前位置。回溯最优前驱得到操作序列；
平局采用预先固定的规则，并保留操作所用的原文坐标。
如“请／与／周五／到场”到“请／于／周五／到场”，只需在零起始区间$[1,2)$
替换“与”为“于”。同样只需一次删除就能把“不得参加”变成“得参加”，
所以编辑距离说明改了哪里，不说明修改是否正确。

假设任务要求输出两句校外参加指南，包含活动日期、地点和精确报名截止条件。
候选甲为“读书会4月18日14:00在海棠楼201举行。校外来宾须在4月16日18:00前报名。”
候选乙为“读书会4月18日14:00在海棠楼201举行。校外来宾无需报名。”
即使教学任务分数分别为$0.78,0.91$，乙也因条件矛盾先被硬筛除，
不能拿高分补偿事实错误。甲进入合格集合后才可与其他合格候选比较简洁程度。
“提交报名信息”若在业务上不等于报名成功，甲还应修订为原通知的准确动作。

数量核验采用另一份明确补充的教学记录：“资料袋共12份”。若甲随后错误生成
“已备21份资料袋”，反馈应保存候选位置、错误值21、记录值12及允许的局部替换，
而不是只说“请检查”。把21改为12后，重新核对整段的实体绑定、日期、数量、句数和格式。
若原始材料存在冲突，不能凭一个数值校验器擅自选值；应返回待澄清状态。
修订达到固定次数仍失败时，交付失败原因或转人工。前章的句对推断和后文的问答核验提供补充证据，
核验器本身也必须用独立标注检查，而不能只让同一生成器认可自己的输出。

\section{数据到文本与开放写作}
\label{sec:nlp-generation-writing}

两项任务都要求写出新句子，内容的确定方式却不同。数据到文本把给定记录表达出来，
未提供的事实通常不允许补写；开放写作可以构造情节，却仍要遵守题目、人物设定和明确的
真实信息边界。内容规划在二者中都能使用：前者规划要说哪些记录，后者还可能决定故事中
将发生什么。

\subsection{结构化数据到文本}
\label{sec:nlp-generation-data-to-text}

输入若已经是一张活动安排表，生成器不必重新识别图片中的表格；它从可用记录出发，
选择词语、合并相关事实、决定指代和句子边界，再形成自然语言。
WebNLG把给定RDF三元组集合与文本配对，为这些表述选择提供任务入口，
但不能据此认为所有报告的宏观内容选择都已经给定
\scite{nlp-intro-gardent2017}

本节用五条规范化记录：$r_1$给出活动名称“机器学习读书会”，$r_2$给出当前时间
“2025年4月18日14:00”，$r_3$给出地点“海棠楼201”，$r_4$给出
“校内师生可直接参加”，$r_5$给出“校外来宾须在2025年4月16日18:00前提交报名信息”。
五条记录均绑定活动编号\code{ML-0418}；$r_4,r_5$的值保留完整的角色与条件。
这是从原通知抽出的当前安排视图，没有原日期字段，因而不能单凭这张表生成
“从4月17日改期”的历史陈述。

\subsubsection{模板与槽位填充}
\label{sec:nlp-generation-template}

当字段和句式比较稳定时，可以把任务规则写成模板：
“\{名称\}于\{当前时间\}在\{地点\}举行。\{校内规则\}，\{校外规则\}。”
模板不需要文本生成训练，参数来自经确认的模式、日期格式和语法规则。
执行时先按活动编号连接字段，统一时区与日期，再填槽；不能把两条不同活动的同名字段
按表格行号直接拼接。

上述五条记录生成两句：
“机器学习读书会于4月18日14:00在海棠楼201举行。
校内师生可直接参加，校外来宾须在4月16日18:00前提交报名信息。”
第一句的三个槽位分别回指$r_1,r_2,r_3$，第二句回指$r_4,r_5$。
若地点缺失，模板应选择明确的缺值分支，如“地点尚未提供”，或者停止生成并报告缺字段；
不能填入默认教室。两条完全相同的地点记录只表达一次，但日期相同而编号不同的活动不能去重。

把同一记录直接序列化为“编号=ML-0418；名称=……；时间=……”也能交给条件生成器。
例如候选“读书会周五下午举行，欢迎参加”语句通顺，却没有表达精确时间、地点和报名规则。
以五条记录为单位核对，它至多完整覆盖$r_1$，不能把模糊的“下午”算作完整覆盖$r_2$；
若反复写两次地点，还应另记冗余。模板的局限是覆盖不了规则之外的句式和组合，
直接生成的局限则包括遗漏和无依据补写，二者都应报告未表达事实与重复事实。

\subsubsection{内容选择与规划后生成}
\label{sec:nlp-generation-content-planning}

字段多到不能全部写入时，需要先决定哪些记录与用途相关，再决定表达次序。
\term{内容选择与规划}（content selection and planning）把这一步显式化：
从记录集合$R$形成记录编号序列$z$，随后按$z$写文本。
Puduppully等的方法以记录编码、选择门和指针式规划产生中间计划，再由带复制的生成器
形成描述；其具体任务为体育记录到报道，下面的活动例子是教学迁用
\scite{nlp-generation-puduppully2019}

有记录、计划和文本三种监督时，可训练
\begin{equation}
  L_{\mathrm{plan}}
  =-\log p_{\bm{\phi}}(z\mid R)
   -\log p_{\bm{\theta}}(y\mid R,z).
  \label{eq:nlp-generation-plan-loss}
\end{equation}
计划中的下一项可以是某条记录的编号或规划结束符。第一项损失监督选择与次序，
第二项监督如何把选定记录说出来。若只有记录与文本，计划并不会自动成为已知标签；
可用记录--文本对齐产生弱监督，或用明确的任务规则构造计划，并记录对齐错误。
对所有可能计划求和的潜变量训练又是另一种计算，不能与上述有计划监督混写。

若用途为“两句校外来宾指南”，五条记录中选择$r_1,r_2,r_3,r_5$，不选$r_4$。
规划为$z=(r_1,r_2,r_3;\,r_5)$，分号表示句子组，得到
“机器学习读书会于4月18日14:00在海棠楼201举行。
校外来宾须在4月16日18:00前提交报名信息。”
选取$r_5$解决的是外宾的参加条件，省略$r_4$是用途允许的取舍；若任务改成完整通知，
同样的遗漏就不合格。反向核验时把每个生成陈述重新对齐到$r_i$，检查计划是否覆盖了要求，
以及文本是否完整实现计划。

名称和时间可通过来源复制保留，混合概率的计算见
第~\ref{sec:nlp-generation-pointer}节。复制能输出“海棠楼201”，却不能阻止模型把它
绑定到另一场活动。与直接生成相比，规划路线增加计划标注、一次规划推断和中间错误来源，
但允许在写句子之前检查选择及顺序。直接生成则一次读取全部记录，可能在内部决定内容，
只是没有可单独验收的显式计划。比较必须把规划消耗的词元与调用也算进总预算。

\subsection{开放写作与续写}
\label{sec:nlp-generation-open-writing}

为小说构造一个新人物不属于事实幻觉，前提是任务明确允许虚构。
本节题目为“失而复得的借书卡”，人物只有林岚，故事发生在同一天的图书馆，
分三段，每段一句，文体平实；允许虚构找卡经过，但不得新增人物替林岚找到卡，
也不得把图书馆业务规则写成真实规定。评价分别检查情节多样性、阅读偏好与约束满足，
不能把任何一个维度叫作全部创作质量。

\subsubsection{Plan-and-Write情节规划}
\label{sec:nlp-generation-plan-write}

直接写第一句时，模型可能还没有决定结尾怎样回应开头。
Plan-and-Write先用有次序的关键词表示故事线，再依据故事线生成文本。
原方法从已有故事每句抽取一个关键词作为规划监督，并比较静态与动态两种方式；
关键词不是完备的事件逻辑，也不是人物状态证明
\scite{nlp-generation-yao2019}

静态方式由题目$u$生成完整计划$z_{1:k}$，再生成故事$y$，
训练分别最大化$p(z\mid u)$与$p(y\mid u,z)$。
动态方式用题目、此前句子和上一规划词预测下一规划词，再据此写下一句；
相应的规划条件是$(u,y_{<j},z_{j-1})$，写作条件还包含新规划词$z_j$。
两个模型可以分别训练并在运行时串接，不应因为流程相连就声称训练一定端到端联合优化。

静态教学计划为“遗失／回找／夹页”。按它形成三段：
\begin{quote}
林岚在图书馆借书时发现借书卡不见了，记起自己刚翻过一本厚书。

她返回阅览桌，把那本厚书从原处取下来，一页页查看。

借书卡正夹在书页之间；她收好卡，回到借书处。
\end{quote}
这里另设一份教学检查记录：第一段林岚在借书处、卡的位置未知；第二段她到阅览桌，
取书后开始找；第三段卡从书页转到她手中，她返回借书处。开头的厚书线索在结尾回收，
人物移动和事件先后可以逐项检查。这份记录是本书的核验手段，不是原模型自动维护的状态表。

动态方式写完第一段后，也可能选择“衣袋”作为第二个规划词，形成
“她先检查外套衣袋，找到夹在票据中的借书卡”，再以“收好”为结尾。
该路线能响应已写内容，但应同步处理第一段“厚书”线索是否仍有作用；
若仍生成“从书中取出刚找到的卡”，便把同一张卡放在两个位置。
静态规划也会失败：计划改为“衣袋”而正文仍沿用“夹页”时，
只有核对最终文本才能发现计划与实现脱节。

\subsubsection{语言模型直接续写与候选重排}
\label{sec:nlp-generation-continuation}

直接续写把题目、人物要求和已有第一段作为前缀，用
第~\ref{sec:nlp-generation-training}节的目标学习后续文本，生成与采样复用
第~\ref{sec:nlp-generation-sampling}节。运行时必须写清是否能看到此前所有段落，
用何种段落标记，以及第三段结束后的EOS；没有EOS而撞上上限的文本仍是截断候选。
多轮共同修改题目或人物设定的对话过程见第~\ref{chap:nlp-question-answering-dialogue}章。

固定上述第一段，教学候选甲续写“返回阅览桌找书，卡夹在书中”；
候选乙续写“管理员把卡送来，林岚感谢他”。
乙新增人物代为找卡，违反已定条件，即便语言偏好评分更高也须淘汰。
在合法候选之间，还可按线索回收、重复和约定风格重排。
重排器的来源可以是人工量表、规则或有标注的偏好模型，来源不同需要分别验证。

比较规划与直接续写时，设两条完整候选、每条最多90个新词元作为共同生成预算。
规划路线每条分配最多12个规划词元和78个正文词元，直接路线每条最多90个正文词元；
两者再接受相同次数的核验，记录实际词元、前向调用和延迟。这个分配只是教学协议，
不是相等运行时间的保证，输入长度与模型结构仍影响成本。
规划在写正文前决定部分情节，直接续写在正文中逐步作决定；完整上下文可用时二者都能
读取已写历史，发生截断时则都可能丢失人物条件。偏好最高的一条仍须检查位置、时间线、
题目要求和允许虚构的边界。

\section{文本编辑}
\label{sec:nlp-generation-editing}

给定一段已经写好的文字，新文本与原文之间就多了一项约定：哪些变化是必要的。
改写调整表达方式，纠错修复错误，简化降低特定读者的理解负担。三者都可能替换词语、
拆句或重写整段，但不能仅凭操作形式判断任务。把“须提交”改成“不必提交”只改了几个字，
造成的信息变化却可能大于一次保持原意的整句重写。

\subsection{改写与风格转换}
\label{sec:nlp-generation-paraphrase}

本章把\term{改写}理解为在约定的意义边界内换一种表达；\term{风格转换}
（style transfer）则另外要求改变正式程度、语气或其他指定属性。
通知变成易读说明时，可以把“现改至”换成“时间已改为”，把参加规则拆成单独一句，
却须保持谁参加、何时参加、谁须提前提交信息。文体不是立场：
“请注意报名截止时间”可以更亲切，“报名不是必须的”则改变了要求。

验收因此分为内容保持、目标风格和表达质量三项。把“校外来宾”换成“全体来宾”
改变适用人群；把14:00改成16:00改变数值；删掉“不得”的“不”改变否定。
即便这些候选更像日常口语，也不能在本任务中接受。

\subsubsection{条件改写与局部编辑}
\label{sec:nlp-generation-conditional-edit}

有成对修订样本$(x,u,y)$时，完整改写以原文$x$和文体要求$u$为条件，
用修订文本$y$训练式~\eqref{eq:nlp-generation-nll}。
例如训练输入“校外来宾须在4月16日18:00前提交报名信息；改为易读说明”，
目标为“校外来宾必须在4月16日18:00前提交报名信息”。
运行时按固定束宽或候选数生成，再检查日期、动作和义务是否保持。
整句生成能改变较大范围的语序，但也须为每处非必要变化负责。

局部路线把同一目标写成有坐标的动作。以
“校外来宾／须／在／4月16日18:00前／提交／报名信息”为原词序列，
零起始区间$[1,2)$执行\code{REPLACE(必须)}，其余位置执行\code{KEEP}。
日期与实体可设为保护片段，禁止修改。多个动作都以本轮原文为坐标，先验证不重叠，
再从左到右重建结果；不能一边插入一边仍按已经移动的位置执行后续动作。
动作标签可以由平行文本对齐得到，也可以来自人工修订日志；两种来源的歧义都应保留。

还可把原文译到另一语言再译回，形成复述候选。例如中间英语
“External guests must submit their registration information before 18:00 on April 16”
回译为“校外来宾必须在4月16日18:00前提交报名信息”，可进入同一个候选集合。
回译只提供不同措辞，不提供等义证书。若回译成“来宾可在4月16日提交信息”，
它同时丢掉人群、义务、精确时点与严格的“前”，应在风格重排之前淘汰。

在三条路线之间固定原文、允许修改范围和候选预算，先用保护片段及语义检查筛选，
再由有文体标注的分类器或人工量表比较正式程度。输入与输出词语高度重合，
仍可能只因一个“不”而反义；文体分类正确，也不能证明角色和数值没有变化。

\subsubsection{Delete--Retrieve--Generate属性转换}
\label{sec:nlp-generation-delete-retrieve}

成对改写难收集时，另一种材料是分别标注属性的句子，而不是逐句对应的修订。
Delete--Retrieve--Generate（\citeyear{li-etal-2018-delete}）据此把局部属性标记与剩余内容分开：
删除原属性片段，从目标属性语料检索适合当前内容的表达，再重组句子。
完整出处见本章文献说明。

属性片段不是预先知道的。设$n_a(g)$为属性$a$语料中短语$g$的出现次数，
可以用带平滑的比值$(n_a(g)+\epsilon)/(n_{\bar a}(g)+\epsilon)$与阈值比较，
把明显偏向$a$的短语作为标记；多属性时分母和对照组另行约定。
删除这些片段得到内容序列$c(x,a)$。检索时在目标属性句子中比较删除后的内容相似度，
以免只找到情感相同、对象却无关的句子。平滑量、阈值、相似度和并列选择规则来自
训练及开发集，不从测试参考倒推。

用明确允许改变评价立场的教学任务说明四种输出方式。原句
“这家书店的服务令人失望”带负面标签，删除“令人失望”后留下
“这家书店的服务”；目标正面语料检索到“这家书店的服务十分周到”。
直接检索基线返回整条检索句，模板基线把“十分周到”放入原句删除位置；
DeleteOnly生成器接收剩余内容和正面标签，可输出“这家书店服务很好”；
DeleteAndRetrieve还接收检索到的目标属性片段，可输出“这家书店的服务十分周到”。
这些是指定输入下的构造候选，不是四个系统的实测结果。若检索句变成
“这家商场的服务十分周到”，整句返回还会把书店改成商场，内容检查便能区分它们。

生成器的监督来自原句重建：以删除后的内容和该句原属性，或其属性片段为输入，
最大化原句概率；训练时对属性片段作扰动可减少机械复制。
运行时再把属性条件换成目标属性。这种训练没有提供真实的跨属性平行目标，
“剩余内容足以重建要保留的信息”仍是需要检验的假设。

对通知则须收窄操作范围。原片段“敬请校外来宾按时提交报名信息”可以删除
“敬请”并检索“请”，生成易读说明；不能因为“须”与正式文体相关就删除义务。
若任务指定保持立场，上面的正负评价转换不是合法比较候选。
同一删除机制可以服务不同任务，但必须重新定义属性之外究竟保留什么。

\subsection{文本纠错}
\label{sec:nlp-generation-correction}

纠错的依据是原文存在错误，不是作者偏好另一种表达。
“请与周五到场”中的同音字、“提高报名流程”中的搭配，以及
“She go to meeting”中的英文词形和冠词，都可以给出明确修订。
“请于周五到场”已经正确，改成“烦请周五莅临”属于文体选择，不能自动算纠错收益。
JFLEG提供英文流畅性修订材料，允许的改动范围比某些最小语法纠错任务更宽，
使用时应沿用其任务定义，而不把它当作中文同音字纠错的直接证据。

三条路线都应保留正确句样本，并区分漏改、正确修改、过改与意义改变。
一轮动作的坐标基于同一原序列；冲突动作先选择兼容集合，重建后才开始下一轮。
输出若将一个正确日期改错，不能用另两处语法修复抵消。多种合法修订则需要多参考、
可接受编辑集合或人工复核，修改级指标见第~\ref{sec:nlp-generation-edit-metrics}节。

\subsubsection{混淆集与语言分数筛选}
\label{sec:nlp-generation-confusion}

若错误多集中在同音、形近或拼写混淆，可以先限制替换范围。
\term{混淆集}是某个位置允许尝试的候选字词集合，它来自词典、输入法混淆统计、
人工规则或有标注错误。集合必须保留原词，否则系统即使遇到正确输入也被迫改写。
给完整候选句计分，例如
\begin{equation}
  s(y;x)=\log p_{\mathrm{LM}}(y)-\lambda\,d(x,y),
  \label{eq:nlp-generation-confusion-score}
\end{equation}
其中语言模型提供上下文分数，$d$为编辑数，$\lambda\geq0$抑制无必要改动。
只有相对原句的得分增量超过开发集确定的阈值，且规则检查通过，才执行替换。

对“请／与／周五／到场”，候选位置取$\{\text{与},\text{于},\text{予}\}$。
设三条完整句子的教学对数分数分别为$-8,-4,-10$，$\lambda=1$，
则保留原文得$-8$，替换“于”得$-5$，替换“予”得$-11$。
阈值为2时，增量3使“于”进入候选，再由时间介词规则核验并恢复
“请于周五到场”。这里数字只用于演示阈值，不代表某语言模型的表现。

领域词反映另一种风险。“本次讨论主题是束搜索”中的“束”虽不如“数”常见，
却是已确认术语；将其保护为词组后，混淆筛选不得输出“数搜索”。
若词典根本没有包含真实误字，搜索不可能找到正确替换。逐位置替换也不能直接补冠词、
删除重复词或调整远处搭配，需要扩充动作或转入下面两条路线。

\subsubsection{GECToR编辑标签纠错}
\label{sec:nlp-generation-gector}

替换词典可扩展为一个有监督的编辑词表。GECToR在各源词位置预测变换标签，
保留大部分原文，仅执行少量修改。原句与正确句的对齐导出
\code{KEEP}、\code{DELETE}、\code{REPLACE\_w}、\code{APPEND\_w}
以及英文词形变化标签，模型对这些标签计算分类交叉熵
\scite{omelianchuk-etal-2020-gector}

标签位置包含句首的特殊位置，以便在第一个词之前插入。
训练预处理用修改后的Levenshtein对齐识别操作；某位置若有多个操作，
一轮只保留首个非KEEP标签，再把剩余修改交给下一轮。
保留正确句参加训练，使全KEEP也是有监督的正常输出。
运行时错误检测阈值和KEEP偏置共同控制是否修改，迭代到没有动作或达到固定轮数。
阈值应在与任务同类的开发集上选择，不把多轮次数当作越多越好。

对“\code{\$START}／She／go／to／meeting”，教学标签为
“KEEP／KEEP／动词第三人称单数变换／APPEND\_the／KEEP”。
同时应用于这一轮原词后，得到“She goes to the meeting”：
变换将go变成goes，在to之后插入the。若目标进一步是
“She goes to the annual meeting”，而每个位置一轮只允许追加一个词，
第一轮仍追加the，第二轮在新序列的the之后追加annual。
第二轮的坐标由第一轮结果重新编号，不能在旧to位置一次执行两个互相覆盖的标签。

编辑词表限制了单轮可表示的变化。未收录的替换、较长短语插入或复杂重排
可能需要多轮，甚至无法在给定预算内完成；候选只改一个词也可能改变否定或日期。
中文必须另定按字还是按词对齐、同音字动作及词语边界，
不能把英文单复数和动词时态标签原样迁入中文。

\subsubsection{条件序列生成纠错}
\label{sec:nlp-generation-seq-correction}

另一条路线直接以错误句为条件生成完整修订句。任务提示可以说明
“仅修正语法，保持实体和数值”，参考修订进入目标损失，EOS参与监督；
用BART或T5等编码--生成模型实现时，预训练结构回指第~\ref{chap:transformer}章，
本节增加的是纠错配对与恢复接口。译码后必须重新对齐原文和结果，才能知道模型改了什么。

固定输入“She go to meeting on Friday”，两项参考修改为go到goes，
以及在to后插入the。生成候选“She goes to the meeting on Monday”
包含三项编辑：前两项正确，Friday到Monday是过改且改变事实。
按这一份明确编辑标注，正确编辑2项、系统编辑3项、参考编辑2项，
精确率为$2/3$，召回率为1；数字的含义与加权汇总见后文。
句子变得更流畅并不能使Monday成立。

与GECToR比较时，可固定最多两轮、每条至多16个新词元或两轮标签前向，
并记录实际调用与运行时间；词元数和标签轮数不是天然相等的计算单位。
整句生成能一次写出“the annual meeting”，而受单次APPEND限制的编辑模型
需要第二轮；编辑模型则保留可见的局部动作，不必重新生成所有正确词。
若正确译法或修订不唯一，例如允许“She attends the meeting on Friday”，
不能仅因它不是最小词形修改就判错，应先确认本任务是否允许这种改写，再扩充参考。

\subsection{文本简化}
\label{sec:nlp-generation-simplification}

简化要回答“让谁更容易理解什么”。读者是否熟悉术语、是否需要完整参加条件、
允许删掉哪些背景，应在生成前确定。ASSET的多参考简化同时包含词汇替换、
拆句和删减等变化，提醒评价不能只奖励缩短。相关任务数据的文献入口见本章文献说明；
真实理解仍需要面向目标读者检查，而不只是自动复杂度下降。

\subsubsection{词典替换、拆句与规则重排}
\label{sec:nlp-generation-rule-simplification}

规则路线从适配读者的词典提出常用表达，再利用句法结构确定拆句位置和省略成分。
例如把“须”解释为“必须”，不把它替换成义务更弱的“可以”；
拆开条件句时，把适用人群带到新句中；代词只有在指代唯一时才用明确名词替代。
词典可由教学词表及人工审订构建，句法与指代分析回用第~\ref{chap:nlp-text-analysis}章。
规则来源是可审查的任务约定，不能称为从少量例子自动发现的通用语言规律。

使用如下带否定与条件的教学长句：
“若您是校外来宾，须在4月16日18:00前提交报名信息；这项要求不适用于校内师生。”
分析得到第一分句条件为“校外”，主干为“须提交”，时间修饰提交动作；
第二分句中“这项要求”指提前提交报名信息。替换“须”、把条件改成显式主语并分句，
可恢复为“校外来宾必须在4月16日18:00前提交报名信息。
校内师生不受这项要求限制。”否定仍只作用于校内师生的要求适用性。

若重排为“请在4月16日18:00前提交报名信息。校内师生不受限制”，
第一句丢掉主语范围，第二句又把一项报名要求扩大成所有限制，两个改动都不合格。
保留原指代比贸然补错指代更安全，必要时应拒绝执行某条规则。
拆句减少了句内依赖，却也可能使跨句指代变难；词频高、句子短或依存树浅
只是可观察属性，不是所有读者理解提升的充分条件。

\subsubsection{ACCESS受控简化}
\label{sec:nlp-generation-access}

普通条件生成把简化幅度隐含在训练资料里，用户难以指定“保留信息但多拆句”。
ACCESS把目标属性编码为输入前缀标记，使同一生成模型能够接收不同简化要求。
它使用字符长度、字符级编辑相似度、词汇复杂度及依存树深度四类代理属性，
按训练句对测量、离散化后附到原句前，再预测简化文本
\scite{martin-etal-2020-controllable}

记原句和目标字符数为$n_x,n_y$，本节长度比采用$n_y/n_x$；
改写程度采用$1-d_{\mathrm{char}}(x,y)/\max(n_x,n_y)$，
其值越大表示文字越相近。词汇复杂度先取各词频排名的对数，
用第三四分位数汇总，再取目标与来源之比。句法复杂度用固定分析器给出的最大依存深度；
本章明确统一为“目标深度／来源深度”，实现时应核对具体工具的比值方向。
字符规范化、词频表、分词器和分析器固定后，属性才可复算；它们并不是真实理解难度。

例如一对教学训练句经固定预处理得到40与36个字符、字符编辑距离12，
词汇对数排名第三四分位数为5与4、最大依存深度为5与3，
四个条件便是$0.90,0.70,0.80,0.60$。
这些是假设的测量输入，不是前述中文句子的实测标注。
按0.05间隔离散并对比值设上限2后，将对应标记与原句拼接，
以真实简化句作目标训练条件最大似然；不能在训练时把待预测目标全文放进输入。
运行时不再测未知目标，而由用户或开发集选定控制值。

对前述带条件的通知，普通生成可产生“校外来宾请提前报名”，
虽短却漏了时点和校内规则；要求保留事实、适度改写与降低句法复杂度的受控候选，
可以是规则例子中的两句完整说明。候选形成后重新计算实际属性，检查是否接近请求值，
同时核对截止时间、义务和人群。控制标记是模型条件而非硬保证；
即使长度比精确达到0.9，也可能删错关键条件。
SARI比较可接受编辑的计算见第~\ref{sec:nlp-generation-edit-metrics}节，
目标读者能否回答“谁须在何时提交什么”则是另一个验收问题。

\section{摘要与信息综合}
\label{sec:nlp-generation-summarization}

摘要不是把原文均匀缩短。日历条目主要保留时间和地点，参加指南必须保留行动条件，
会议纪要还须分清建议、已作决定和待办事项。单来源摘要从一份材料选择和组织内容；
多来源综合另外要处理转载、互补与冲突。两者都可能直接取句、融合片段或生成新句，
但每个事实陈述仍应能找到相应来源。

\subsection{单来源摘要}
\label{sec:nlp-generation-single-summary}

通用摘要由任务预先规定重要信息，查询导向摘要则以用户关注点改变选择。
对共同通知，通用版本可以报告改期、地点和两类参加规则；
查询“校外来宾何时、到哪里、怎样参加”时，校内规则可以省略，
校外截止条件却不可省略。XSum所研究的极短摘要强调高度压缩的概括，
不能据此把所有通知压成一句而不管用途。相关基准的出处统一列于本章文献说明。

对候选应分别检查选择是否回应关注点、信息是否重复、句间组织是否清楚、
所说内容是否受来源支持，以及关键内容是否遗漏。抽取原文不自动通过这些检查：
单独摘出“这项要求不适用”却不带上人群或指代，会把正确原句变成误导性的摘要。

\subsubsection{Lead与TextRank抽取摘要}
\label{sec:nlp-generation-extractive}

Lead直接按原文顺序取前若干句，直到句数或词元预算耗尽。
它不需要摘要训练样本，适合用作可复现基线，但依赖重要信息是否位于开头。
TextRank改用句子间联系估计重要性：分句后把各句作为节点，
用词重合或其他明示相似度赋边权，再迭代分配分数。
图排序原理已在前章介绍，这里只展开句子选择和恢复
\scite{nlp-mihalcea2004textrank}

设$w_{i,j}\geq0$为不同句子间的对称边权，令自环权为零，
本节把分数归一化为总和1：
\begin{equation}
  r_i^{(k+1)}
  =\frac{1-d}{n}
   +d\sum_{j:\,w_{j,i}>0}
   \frac{w_{j,i}}{\sum_\ell w_{j,\ell}}r_j^{(k)}.
  \label{eq:nlp-generation-textrank}
\end{equation}
$d$是沿图传播的权重；没有出边的句子按均匀分布分配其质量。
以均匀分数启动，迭代到相邻两轮差小于预定阈值或达到上限。
按分数选句不等于按分数排列摘要：选完通常恢复原次序，以免打乱条件和指代。

将通知整理为三句教学输入：
$s_1$“读书会改至4月18日14:00举行”，
$s_2$“读书会4月18日在海棠楼201举行”，
$s_3$“校外来宾参加读书会须在4月16日18:00前提交报名信息”。
为能手算，明确只保留如下检索词集合：
$K_1=\{\text{读书会},\text{4月18日},\text{14:00}\}$，
$K_2=\{\text{读书会},\text{4月18日},\text{海棠楼201}\}$，
$K_3=\{\text{读书会},\text{校外来宾},\text{报名信息}\}$。
取$w_{i,j}=|K_i\cap K_j|$，边权为$w_{1,2}=2,w_{1,3}=w_{2,3}=1$。
这些词集合只用于排序，恢复时仍取完整句，不能把截止时间从正文删掉。

令$d=0.85$，从初值$(1/3,1/3,1/3)$开始迭代：
\[
  \begin{aligned}
    \text{第一轮}&:\quad(0.38056,0.38056,0.23889),\\
    \text{第二轮}&:\quad(0.36718,0.36718,0.26565),\\
    \text{收敛值}&:\quad(0.37013,0.37013,0.25974).
  \end{aligned}
\]
预算为两句、平局取靠前句时，
Lead和该TextRank都选择$s_1,s_2$，重复了日期却遗漏校外要求。
若查询是“校外来宾如何办理”，显式查询相关分数可以把$s_3$排到首位，
再选择所需时间或地点。也可把均匀跳转换为查询先验，但其强度仍须验证。
选后的去重见第~\ref{sec:nlp-generation-mmr}节；高中心性本身不保证回答用户的问题。

\subsubsection{pointer-generator复制与生成}
\label{sec:nlp-generation-pointer}

抽取整句难以合并日期和地点，纯词表生成又可能写错罕见名称。
\term{指针生成网络}（pointer-generator）把两种输出分布混合：
一部分概率从固定词表生成，另一部分概率由来源注意分配到源词。
同一个词在来源出现多次时，其位置概率要相加，而不是只取最高一次
\scite{nlp-intro-see2017}

设解码步$t$对来源位置$i$的注意为$a_{t,i}$，总和为1；
生成门$g_t\in[0,1]$由解码状态、来源上下文和当前输入经可学习映射得到。
词表分布为$p_{\mathrm{vocab},t}$，则扩展词表上的概率为
\begin{equation}
  p_t(w)=g_t p_{\mathrm{vocab},t}(w)
  +(1-g_t)\sum_{i:x_i=w}a_{t,i}.
  \label{eq:nlp-generation-pointer}
\end{equation}
不在固定词表的词令第一项为零；固定词表中但来源没有的词令第二项为零。
扩展词表为来源中的未登录词分配临时编号，解码后按该样本的映射恢复原字符串。
训练目标仍是参考词的负对数概率，门与注意随目标损失学习，并非人工逐步指定复制标签。

教学来源词元为“读书会／地点／海棠楼201／请／到／海棠楼201”，
注意分布为$(0.05,0.05,0.45,0.10,0.05,0.30)$，生成门$g_t=0.4$。
若“海棠楼201”是未登录词，它的输出概率为
$0.6(0.45+0.30)=0.45$；若“举行”的词表概率为0.5且来源未出现，
其总概率为$0.4\times0.5=0.20$。
以名称为参考时，这一步损失为$-\log0.45\approx0.7985$。
若只取一个名称位置，就会错误地把名称概率算成0.27。

为减少反复关注已经概括过的片段，coverage维护此前注意的累积
$c_{t,i}=\sum_{t'<t}a_{t',i}$，把它作为注意计算的额外输入，
并增加$\sum_i\min(a_{t,i},c_{t,i})$的重叠惩罚。
假设此前一步注意为$(0,0,0.6,0.1,0,0.3)$，
本步重叠为$0.45+0.10+0.30=0.85$；更新后的累计注意总和为2。
这不是强迫模型把所有源词各读一次，也不是禁止重复词：
摘要仍可能需要重复活动名，而来源也有不应选入的背景。
复制了正确名称却把它接到错误活动上，或重复强调一个无关片段，都仍需任务核验。

\subsubsection{BART与T5的条件摘要}
\label{sec:nlp-generation-conditional-summary}

预训练编码--生成模型可以把整份材料编码为条件，再生成目标摘要。
BART接入时以来源及明确的任务条件为输入，以摘要作输出；
T5也可以用任务前缀把“摘要”及其要求写入文本接口。
二者的预训练目标不同，摘要训练则都可采用式~\eqref{eq:nlp-generation-nll}，
不能仅凭模型名称推断系统是否读取了查询、是否微调以及能接收多长输入。
T5的通用文本接口参见\citetitle{raffel2020t5}；
BART及本节相关任务文献的完整条目见章末文献说明。

仍以“校外来宾何时、到哪里、怎样参加”为查询，单份通知按原位置拆成五个事实片段：
原定日期、当前日期与时刻、地点、校内规则、校外截止条件。
保留每个片段的字符范围和通知编号。查询相关性筛选保留当前时间、地点和校外条件；
若分句导致两个片段都重复报告当前日期，再用MMR扣除已经选择的重复信息，
而不是把相关性前几名直接拼接。MMR的逐轮计算在下一小节给出。
选后按“活动安排；参加前提”组织成两组，条件生成器读取两组及其来源位置，
输出“读书会于4月18日14:00在海棠楼201举行。
校外来宾须在4月16日18:00前提交报名信息。”

核验把第一句拆回日期、时刻与地点，分别对应通知的当前时间片段和地点片段；
第二句的角色、动作和截止时点均对应校外规则。原定日期及校内规则未表达，
是查询允许的选择，不计为本次关键遗漏。若查询改为“这次安排有哪些变更”，
原定日期就须重新进入候选。引用整份通知只能说明来源范围，不能替代这一步陈述对齐。

直接输入全文省去外部选择，生成器仍需在内部决定讲什么；
选择后生成降低输入长度并暴露中间依据，却有不可恢复的筛选错误。
若选择器误删了校外条件，生成器再强也不能从留下的片段可靠地重建。
比较时应把选择、重排、生成和核验的总调用与输入词元都计入成本。

长文还存在上下文长度限制。分段摘要再合并时，各段保留
“陈述、原文范围、是否含条件或待决状态”，合并阶段沿引用回查原文。
例如会议材料先说“林岚建议改至周五”，后面才说“与会者否决该建议”，
若第一阶段只保存“会议改至周五”，第二阶段已经失去恢复建议状态的依据。
直接长上下文可同时读取前后发言，但并不保证模型正确利用它们。
两条路线都要检查最终决定，并对行动项恢复“负责人、动作、截止时间、决定依据”。
教学会议若明确决定“由林岚在4月16日18:00前汇总报名信息”，
纪要可以列出这一行动项；原通知没有该负责人，不能把会议补充事实反写成通知已说过。

\subsection{多来源信息综合}
\label{sec:nlp-generation-multi-summary}

来源增多后，重复不能当作多份独立证据，表面不一致也不一定是矛盾。
旧通知的4月17日和改期通知的4月18日对应不同版本；
另一活动也在4月18日举行，则属于不同对象。应先对齐
“对象、属性、值、适用时点、来源”，再判断是转载、互补、条件差异还是冲突。
确认版本关系后可表达最新状态；两个同样有效而相互冲突的来源应并列表述或请求澄清，
不能按转载数量投票。
图~\ref{fig:nlp-generation-provenance}展示同一活动的三份教学通知中的重复、互补和版本差异。

\begin{figure*}[htbp]
  \centering
  % Adapted from academic-drawing/templates/tikz_template.tex to the book palette.
% Source fragments are grouped by claim, not by chronological reading order.
\begin{tikzpicture}[
  every node/.style={font=\figurefont\small,align=center,line width=.6pt},
  source/.style={book node,text width=59mm,minimum height=12mm,
    fill=FigureInput!10,draw=FigureInput,font=\figurefont\small},
  claim/.style={book node,text width=65mm,minimum height=14mm,
    fill=FigureOutput!10,draw=FigureOutput,font=\figurefont\small},
  support/.style={book arrow,rounded corners=1mm},
  qualifier/.style={support,dashed,draw=BookAmber}
]
  \node[font=\figurefont\bfseries] at (-47mm,10mm) {来源片段};
  \node[font=\figurefont\bfseries] at (45mm,10mm) {摘要陈述及保留状态};
  \node[source] (date1) at (-47mm,0)
    {$D_1$：4月15日官方改期通知\\4月18日14:00};
  \node[source] (date2) at (-47mm,-15mm)
    {$D_2$：同日转载\\4月18日14:00};
  \node[claim,minimum height=23mm] (date) at (45mm,-7.5mm)
    {当前时间：4月18日14:00\\$D_1,D_2$支持同一陈述\\转载去重，不增加独立证据};
  \node[source] (place) at (-47mm,-32mm)
    {$D_2$：地点片段\\海棠楼201};
  \node[claim] (placeout) at (45mm,-32mm)
    {地点：海棠楼201\\与日期互补};
  \node[source,minimum height=19mm] (rule) at (-47mm,-51mm)
    {$D_1$：校外参加条件\\须在4月16日18:00前\\提交报名信息};
  \node[claim,minimum height=19mm] (ruleout) at (45mm,-51mm)
    {校外来宾须在4月16日18:00前\\提交报名信息\\保留角色、义务与截止条件};
  \node[source,draw=BookAmber,fill=BookAmber!8] (old) at (-47mm,-70mm)
    {$D_3$：4月14日旧通知，现已归档\\原定4月17日};
  \node[claim,draw=BookAmber,fill=BookAmber!8] (oldout) at (45mm,-70mm)
    {旧版日期：4月17日\\不作为当前安排};
  \node[source] (inside) at (-47mm,-86mm)
    {$D_3$：校内参加规则\\校内师生可直接参加};
  \node[claim,draw=BookMuted,fill=BookPaper] (omission) at (45mm,-86mm)
    {本次校外指南可省略\\仍保留原文，供其他用途调用};
  \node[draw=FigureLoss,fill=FigureLoss!7,rounded corners=1mm,
    text width=150mm,minimum height=11mm] at (-1mm,-103mm)
    {无支持新增：“主讲人为某教授”。三份资料均未给出主讲人，不能写入事实摘要。};
  \draw[support] (date1.east) -- (-3mm,0) |- ([yshift=4mm]date.west);
  \draw[support] (date2.east) -- (-3mm,-15mm) |- ([yshift=-4mm]date.west);
  \draw[support] (place.east) -- (placeout.west);
  \draw[support] (rule.east) -- (ruleout.west);
  \draw[qualifier] (old.east) -- node[above,font=\figurefont\footnotesize] {旧版} (oldout.west);
  \draw[support,draw=BookMuted] (inside.east) -- (omission.west);
\end{tikzpicture}

  \caption{摘要陈述与来源的对应。三份教学通知按活动编号\code{ML-0418}对齐；
  实线表示支持，虚线表示需要版本解释的日期差异。转载的相同日期合为一个事实，
  地点与报名条件分别保留来源。灰框标明用途允许省略的校内规则，
  底部警示框标明无来源支持的新增陈述，二者不能混为一类。}
  \label{fig:nlp-generation-provenance}
\end{figure*}
\FloatBarrier

图中$D_1$为4月15日官方改期通知，含当前时间和校外规则；
$D_2$为同日转载，重复当前时间并补列地点；
$D_3$为4月14日发布、现已归档的旧通知，给出原定4月17日及校内规则。
它们都是本节构造的来源片段，不是实际检索结果。
同一陈述可以有多个来源，但归档日期不能与当前日期一起无条件拼接。
SummEval提供摘要质量维度的评价入口，FActScore提供原子事实支持率的入口；
事实支持率不度量应写信息是否全部写出，具体核验见
第~\ref{sec:nlp-generation-faithfulness}节和第~\ref{chap:nlp-question-answering-dialogue}章。

\subsubsection{MMR与最大覆盖的选择式综合}
\label{sec:nlp-generation-mmr}

只按相关性排序，容易让同一条重要日期占满摘要。
\term{最大边际相关性}（maximal marginal relevance, MMR）每次既考虑当前片段
对查询的相关性，也扣除它与已经选入内容的最大相似度
\scite{nlp-generation-carbonell1998}
对已选集合$S$和尚未选入的合法片段$i$，选择准则为
\begin{equation}
  i^*=\argmax_{i\notin S,\;\operatorname{cost}(S\cup\{i\})\leq B}
  \left[\lambda\,\operatorname{rel}(i,q)
  -(1-\lambda)\max_{j\in S}\operatorname{sim}(i,j)\right].
  \label{eq:nlp-generation-mmr}
\end{equation}
其中$q$为查询，$B$为预算，$\lambda\in[0,1]$控制相关性与去冗余的权衡；
$S$为空时把最大相似度定义为零。本节每轮取可容纳且分数为正的最高项，
直到预算不足、没有正分项或候选耗尽。相似度和相关性应采用可比较的标度，
阈值及$\lambda$在开发集确定。

在三份资料中，先把$D_1,D_2$逐字重复的当前时间聚成一个片段，
保留来源集合$\{D_1,D_2\}$；确认$D_3$为旧版本后，将旧日期保留在历史记录中，
不把它列为当前安排候选。此处近重复聚合是独立前处理，不是MMR自动判断来源独立性。
另一个较短的日期重述仍可能通过聚合阈值，因此需要选择阶段继续去冗余。
教学候选与固定分词成本如下，$D_2$的重述是本节显式补充的同义片段。

\begin{fullwidth}
\begin{center}
  \begin{tabular}{@{}clrr@{}}
    \toprule
    片段 & 按斜线分词的内容 & 成本 & 相关性\\
    \midrule
    $A$ & 读书会／改至／4月18日／14:00 & 4 & 0.95\\
    $B$ & 读书会／时间／已更新／至／4月18日 & 5 & 0.92\\
    $C$ & 地点／为／海棠楼201 & 3 & 0.80\\
    $D$ & 校外来宾／须／在／4月16日18:00／前／提交／报名信息 & 7 & 0.85\\
    \bottomrule
  \end{tabular}
\end{center}
\end{fullwidth}
表中成本只用于教学选择预算，实际正文分词器不必按这些词块切分。
$A,B$相似度为0.90，$A,C$为0.10，$A,D$为0.20，
$B,C$为0.10，$B,D$为0.15，$C,D$为0.10。
所有分数均为给定的教学输入，不是向量模型测量。
$A$来自$D_1,D_2$，$B,C$来自$D_2$，$D$来自$D_1$。

取$\lambda=0.7,B=14$。首轮分数为
$A:0.665,B:0.644,C:0.560,D:0.595$，选$A$。
第二轮$B$减去$0.3\times0.9$后为0.374，
$C$为0.530，$D$为0.535，故选$D$；已用成本11。
第三轮只有$C$仍可容纳，其最大已选相似度为0.10，得分0.530，
选入后成本恰为14。若只按相关性从高到低扫描可容纳项，
会选$A,B$，剩余5无法容纳$D$，继而选$C$，遗漏校外条件。
这个比较说明当前构造中的预算取舍，不证明任意MMR权重都优于相关性排序。

加权最大覆盖换一个角度计分：不直接惩罚句子相似，而是奖励新覆盖的事实。
设日期、时刻、地点、完整校外规则四个事实权重分别为$2,1,2,4$；
$A$覆盖前两项，$B$只覆盖日期，$C,D$分别覆盖后两项。
若$c_i$是片段成本，目标是在$\sum_{i\in S}c_i\leq14$下最大化
$\sum_f w_f\mathbb{I}[f\text{被}S\text{覆盖}]$。
以“新增权重／成本”作教学贪心，初值为
$A:3/4,B:2/5,C:2/3,D:4/7$，先选$A$；
此后$B$新增权重为0，再选$C,D$，覆盖权重9，成本14。
相关性基线的$A,B,C$仅覆盖权重5。一般的带预算覆盖问题不能因这个小例子
就宣称贪心总能取得最优解，事实抽取本身也会出错。

选取次序$A,D,C$不是阅读次序。把$A,C$组合为“何时何地”，
再表述$D$的参加前提，得到上一小节的两句摘要。
按陈述保存$A\leftarrow\{D_1,D_2\}$、$C\leftarrow D_2$、$D\leftarrow D_1$，
核验日期与时刻共同来自$A$，而截止时间不能从$D$误接成活动时间。
若$D_3$不是归档旧版而是同日有效通知，冲突尚未解决，就不能在聚合时悄悄删除它；
应另保留冲突标记，并让最终表达说明不确定性。

\subsubsection{词图最短路径压缩与多句融合}
\label{sec:nlp-generation-word-graph}

选择式综合仍保留原句的冗长部分。对已经确认表达同一事件的近重复句簇，
可以把词语及相邻关系构成有向图，再寻找一条较短而受来源支持的表达路径。
词形和词性相同的词优先共用节点，有多个可能位置时依据相邻上下文区分；
起止节点连接每个句子的首尾。图可能有环，不能把它默认当作有向无环图
\scite{filippova-2010-multi}

原工作比较过不同边权与重排方式。为便于复核，本节采用其基础路线的倒数频率权重：
一条边在句簇出现$f(u,v)$次，成本为$1/f(u,v)$。
枚举若干条无环的短路径，以是否包含动词、是否达到最小长度等条件过滤，
再按预定分数选择。下面的最小长度设为五个词元，是教学条件，
不是把某篇论文的英语开发集长度阈值照搬给中文。

输入三句，分词固定为
\begin{quote}
读书会／于／周五／在／海棠楼201／举行。\\
本次／读书会／周五／在／海棠楼201／举行。\\
读书会／将／于／周五／在／海棠楼201／举行。
\end{quote}
三句同指已确认的当前活动，不包含报名条件冲突。
“起点到读书会”出现2次，成本$1/2$；
“于到周五”出现2次，成本$1/2$；
共同尾部“周五到在、在到海棠楼201、海棠楼201到举行、举行到终点”
各出现3次，各成本$1/3$；其他边各出现1次，成本1。

路径“起点、读书会、周五、在、海棠楼201、举行、终点”的成本为
$1/2+1+4/3=17/6$，恢复为“读书会周五在海棠楼201举行”。
它使用第一、三句的直接起点边和第二句的“读书会到周五”边，
因此并非简单取某一完整原句。经“于”的路径成本为
$1/2+1+1/2+4/3=10/3$，经“本次”的对应路径成本为$13/3$。
最短候选有五个正文词元且包含动词“举行”，通过这两项形式过滤。
随后还要核对“周五”对应4月18日、地点与动作属于同一活动。

直接选第一句会保留“于”，词图路线在本例删除了冗余前缀与虚词，
但不能据此把另一簇“校外来宾须提前提交”随意接入。
若把不同角色的句子错误聚为一簇，合法路径也可能生成“校内师生须提前提交”。
连通性、低成本和含动词只检查图与表面形式，不证明关系及语法正确；
需要融合互补事实时，下面的内容规划更便于显式保持其条件。

\subsubsection{内容规划与多文档条件生成}
\label{sec:nlp-generation-multi-plan}

多文档规划将事实单元保存为
“活动编号、事件或属性、值、有效时点、适用条件、来源位置”，
再按主题或时间把它们排成句子组。记录到文本的训练与规划机制
复用第~\ref{sec:nlp-generation-content-planning}节，
此处增加的是跨来源身份对齐和每个输出陈述的证据恢复。
来源标识不能只在生成结束时由模型随意补写。

另加一条明确的教学记录：编号\code{IR-0418}的“检索系统研讨会”
也在4月18日14:00举行，地点为松柏楼305。若错误内容表只以日期作连接键，
就可能把读书会名称接到松柏楼305，再生成一条语法完全正确的错误通知。
修正后分别保存
$(\code{ML-0418},\text{地点},\text{海棠楼201},D_2)$与
$(\code{IR-0418},\text{地点},\text{松柏楼305},D_4)$，
其中$D_4$只是这条新增记录的来源名，不改变图中的三份通知。
按对象分组可恢复两句综合：
“机器学习读书会于4月18日14:00在海棠楼201举行。
同日14:00，检索系统研讨会在松柏楼305举行。”
若用途仍是校外参加指南，读书会的报名条件还须另列，不能因固定两句展示而判为完整。

全文直接生成同时读取各来源，没有显式内容表，可能隐含地完成绑定，也可能误合并。
规划后生成增加抽取与规划调用，却能在表达之前检查连接键、版本和冲突。
分段后合并降低单次输入长度，但第一阶段若删掉活动编号，
第二阶段就无法仅凭同日同地可靠地恢复对象。三条路线应在同样来源、总输入输出预算
和核验调用数下记录实际成本，并逐项检查事实覆盖与引用支持。
引用一份确实含有“松柏楼305”的资料，仍不足以支持“读书会在松柏楼305”；
支持对象是完整关系，而不是孤立词语。

\section{机器翻译}
\label{sec:nlp-generation-translation}

翻译要把原意传到另一种语言，允许语序调整和句子拆合，不要求每个词一一对应。
句子翻译主要依据当前句及明确背景，篇章翻译还处理跨句指代、术语和论述关系。
低资源、领域、实时性及交互修订改变可用信息与预算，不另行改变忠实传达原意的要求。
输入或输出为声音时，声学识别与合成由第~\ref{chap:nlp-speech-tasks}章展开。

\subsection{句子翻译}
\label{sec:nlp-generation-sentence-translation}

固定语言方向、领域和允许背景后，译文要分别核对实体、数量、否定及关系。
共同通知中“校外来宾”不能只译成guests，“4月16日18:00前”
不能只译成on April 16；“提交报名信息”也不必等于已经取得参加资格。
NLLB和WMT24分别提供多语及跨领域翻译的研究入口，出处见本章文献说明。
一个系统支持某种语言，并不表示已经在当前领域、语向和用途中达到所需质量。

\subsubsection{IBM词对齐与短语统计翻译}
\label{sec:nlp-generation-ibm}

平行句子通常没有标明哪个源词对应哪个目标词。IBM Model 1把对应关系作为隐变量，
用词翻译概率解释观察到的目标句，再依据当前模型推断对应关系并更新概率。
本节固定源、目标句长，避免把句长模型与词对齐计算混在一起
\scite{brown-etal-1993-mathematics}

设源句为$x_{1:n}$，另外增加$x_0=\mathrm{NULL}$，表示目标词没有显式源词对应。
目标句为$y_{1:m}$，隐变量$a_j\in\{0,\ldots,n\}$表示$y_j$对齐的位置，
$t(v\mid u)$为源词$u$生成目标词$v$的概率，满足$\sum_v t(v\mid u)=1$。
均匀独立的对齐先验给出
\begin{align}
  p(y,\bm{a}\mid x,m)
    &=\frac{1}{(n+1)^m}\prod_{j=1}^{m}t(y_j\mid x_{a_j}),\\
  p(y\mid x,m)
    &=\frac{1}{(n+1)^m}\prod_{j=1}^{m}
       \sum_{i=0}^{n}t(y_j\mid x_i).
  \label{eq:nlp-generation-ibm-likelihood}
\end{align}
第二式对所有对齐求和；由于每个目标位置独立选择源位置，求和可以逐位置分解。
这也暴露了限制：在给定词袋下改变源词次序，不会改变Model 1的句子概率。

EM的E步用当前$t$计算对齐后验
\begin{equation}
  \gamma_{j,i}
  =p(a_j=i\mid x,y,m)
  =\frac{t(y_j\mid x_i)}{\sum_{\ell=0}^{n}t(y_j\mid x_\ell)}.
  \label{eq:nlp-generation-ibm-posterior}
\end{equation}
同一个源词重复出现时，每个位置都进入分母和计数。
把全语料中源词$u$与目标词$v$的后验相加，得到期望计数$c(u,v)$。
M步在每个源词下归一化，
$t_{\mathrm{new}}(v\mid u)=c(u,v)/\sum_{v'}c(u,v')$；
分母为零的未观测源词需另定保留或平滑规则，不能直接除零。

\Needspace{12\baselineskip}
\begin{example}[含空词的一轮EM]
  \label{ex:nlp-generation-em}
  教学语料只有两对句子：第一对为“我／走”与“I／go”；
  第二对为“我”与“I”。
  目标词表仅有I和go，初始概率及累计期望计数如下：
  \begin{center}
    \begin{tabular}{@{}lrrrr@{}}
      \toprule
      源词 & $t(\mathrm{I}\mid u)$ & $t(\mathrm{go}\mid u)$
        & $c(u,\mathrm{I})$ & $c(u,\mathrm{go})$\\
      \midrule
      NULL & $1/2$ & $1/2$ & $28/39$ & $1/3$\\
      我 & $4/5$ & $1/5$ & $224/195$ & $2/15$\\
      走 & $1/5$ & $4/5$ & $2/15$ & $8/15$\\
      \bottomrule
    \end{tabular}
  \end{center}
  第一对中，I的分母为$1/2+4/5+1/5=3/2$，
  在NULL、我、走上的后验为$(1/3,8/15,2/15)$；
  go的分母也为$3/2$，后验为$(1/3,2/15,8/15)$。
  第二对只有NULL和我，I的分母为$1/2+4/5=13/10$，
  后验为$(5/13,8/13)$。因此NULL到I累计
  $1/3+5/13=28/39$，我到I累计$8/15+8/13=224/195$，
  其余计数直接来自第一对。

  三个源词的计数总量分别为$41/39,50/39,2/3$，相加为3，
  正好等于全语料三个目标位置。M步得到
  \[
    \begin{array}{c|cc}
      &t_{\mathrm{new}}(\mathrm{I}\mid u)&t_{\mathrm{new}}(\mathrm{go}\mid u)\\
      \hline
      \mathrm{NULL}&28/41&13/41\\
      \text{我}&112/125&13/125\\
      \text{走}&1/5&4/5
    \end{array}
  \]
  “我”生成I的概率从0.8增加到0.896，因为第二对额外提供了对应证据。
  固定句长的语料似然从$0.25\times0.65=0.1625$变为约0.19054。
  第一对各目标词取最高后验时，I对齐“我”，go对齐“走”；
  NULL的计数仍保留，不能因这次硬对齐未选到NULL就把它从EM中删除。
\end{example}

词对齐还不是整句翻译。短语统计翻译在对齐上提取相互一致的连续源、目标词块：
短语内部词的对齐不能越过另一侧选定边界，未对齐边界词的扩展另按规则处理。
这里“短语”指连续词块，不必是句法短语。
在独立教学句对“我／今天／出发”与“I／leave／today”中，
给定对齐我--I、今天--today、出发--leave，
源块“今天／出发”可对应目标块“leave／today”。
它在一个词块内表达了语序变化，逐词独立替换本身没有这种次序信息。

运行时把源句分成已覆盖词块，选择目标词块并决定其次序，
结合短语翻译、目标语言模型、重排和长度等特征计分搜索。
特征权重用开发集调节，短语概率由平行语料计数估计。
对同一源句，教学候选“I today leave”的三项加权分数为
$-0.5,-3.0,-0.1$，合计$-3.6$；
“I leave today”为$-0.8,-1.0,-0.4$，合计$-2.2$。
后者虽然短语翻译项较低，综合目标语言顺序后仍胜出。
这只是明确特征下的搜索示例；短语抽取和模型出处见章末文献说明。
对齐是否准确、源词是否覆盖完整以及译文是否传达原意，需要分别评价。

\subsubsection{神经条件翻译与回译训练}
\label{sec:nlp-generation-neural-translation}

神经路线直接从平行样本学习$p(y\mid x,u)$，其中$u$标明语言方向及允许的领域条件。
源句按固定分词器编码，目标端以BOS启动、移位读取参考并预测到EOS，
损失和束搜索复用第~\ref{sec:nlp-generation-training}、
第~\ref{sec:nlp-generation-beam}节。多语共享模型可把目标语言标记放入条件，
但语言标记不自动提供领域术语；术语表、检索上下文及其来源须另外记录。

例如输入共同通知的校外规则，目标词元恢复后应表达
“External guests must submit their registration information before 18:00 on April 16.”
检查external guests保留角色，must保留义务，before保留严格截止关系，
日期、时刻和提交对象均有来源。另用“不要修改12份记录”的教学句检查否定与数量；
“Modify the 21 records”即使语法正确，也同时犯了否定和数值错误。
目标语言词元合并后的标点与数字检查不能只在子词级进行。

平行语料不足而目标语单语较多时，可以用\term{回译}（back-translation）构造额外配对。
先用反向模型$q(x\mid y)$把目标语单语$y$译成合成源文$\widetilde{x}$，
再将$(\widetilde{x},y)$与真实平行样本混合训练正向$p(y\mid x)$。
例如英文单语“Do not modify the 12 records”经过反向模型得到
“不要修改12份记录”，正向训练仍以原英文为目标。
真实数据与合成数据可用明确的采样比例混合；若写为损失，
可取两个数据集各自平均NLL的加权和，权重在开发集确定
\scite{sennrich-etal-2016-improving}

回译中的目标文本来自单语资料，不等于合成源文一定准确。
若反向模型漏掉“不”，这对训练样本就把肯定源句和否定目标配在一起，
应通过质量筛选、降权或人工检查控制这种噪声。筛选器不能偷看独立测试参考。
从平行文本之外取得额外语言材料，与改变正向网络结构是两项不同变化。

\term{枢轴翻译}是在运行时经过第三种语言，并非上述数据扩充。
例如将“不要修改12份记录”先译成英语“Do not modify the 12 records”，
再译成德语“Ändern Sie die 12 Datensätze nicht.”，最终只交付德语。
两步各有解码及恢复成本；若第一步已经漏掉not，
第二步仅凭肯定的英语中间句没有依据恢复否定。
回译噪声通过训练样本影响参数，枢轴错误则沿两步推理传播，二者应分开诊断。

\subsubsection{STACL与wait-k前缀翻译}
\label{sec:nlp-generation-wait-k}

实时翻译不能总等源句完整到达，但过早生成可能遇到尚未出现的动词或限制词。
STACL把训练和解码都限制在已读源前缀上；
wait-k先读$k$个源词，此后通常读一个、写一个。
第$t$个目标词允许看到的源长度为
\begin{equation}
  g(t)=\min(k+t-1,n),\qquad
  p(y\mid x)=\prod_t p(y_t\mid x_{1:g(t)},y_{<t}).
  \label{eq:nlp-generation-wait-k}
\end{equation}
训练也采用这个可见性限制，不能让双向源编码偷偷读到后缀后再只遮住解码注意。
原方法在源文未读完时按前缀贪心输出，源文结束后可对尾部采用束搜索。
方法出处为\citetitle{nlp-intro-ma2019stacl}（\citeyear{nlp-intro-ma2019stacl}），
完整条目见章末文献说明。

用日语教学源句“彼は／会議に／参加／し／ない”说明句末否定。
五个源词块依次表示“他、会议、参加、做、否定”，完整句意是
“He will not attend the meeting”。分词块仅为展示READ和WRITE轨迹。
设$k=2$，每步只保留一个后继，最多生成八个目标词元含EOS，
构造如下候选过程：
\begin{center}
  \begin{tabular}{@{}clll@{}}
    \toprule
    目标步 & 新READ内容 & 已读长度 & WRITE\\
    \midrule
    1 & 彼は、会議に & 2 & He\\
    2 & 参加 & 3 & will\\
    3 & し & 4 & attend\\
    4 & ない，并收到源结束信号 & 5 & the\\
    5 & 无 & 5 & meeting\\
    6 & 无 & 5 & EOS\\
    \bottomrule
  \end{tabular}
\end{center}
第三步选择attend时还没有看到否定；假定它在当时分布中高于not，
贪心就作出了过早的肯定承诺。第四步虽读到ない，也不能在只追加词元的接口中
把not插回will与attend之间；该完成候选漏译否定，必须判为失败。
这不是说任何短等待策略都必然失败，而是展示一个有完整可见性限制的失效轨迹。

若同样最多八步、束宽一，改为$k=5$，第一次WRITE前即可读取完整源句，
候选“He／will／not／attend／the／meeting／EOS”可以在七步完成。
首次输出前分别等待2与5个源词，质量、实际词元和等待不同。
源结束后继续生成直至目标EOS或预算耗尽，不能把源句结束直接当作目标结束。
真实延迟还取决于读写速率、模型耗时和输入粒度，不能只报告$k$。
允许用户或系统撤回早期译文时，可以把attend前插入not作为一次修订，
但那改变了交付接口，应额外报告撤回范围和修改代价；
声音链路的等待见第~\ref{chap:nlp-speech-tasks}章。

\subsection{篇章翻译}
\label{sec:nlp-generation-document-translation}

逐句看似正确的译文放在一起，仍可能指代不明或术语不一致。
组织名重复出现时是否指同一机构，省略主语是否延续前句对象，
“然而”是否保留论证转折，都需要超出当前句的信息。
篇章质量还依赖用途：读者能否据译文按要求行动，与各句是否自然不是同一问题。
DiscoX及MQM提供专业篇章和错误评价入口，但应按具体语向、领域和标注协议解释其结果。

\subsubsection{上下文条件神经翻译}
\label{sec:nlp-generation-context-translation}

上下文条件翻译在当前源句之外加入前后句、检索到的相关段落或全文，
并明确标出当前应翻译的片段。文档级样本必须保留句子顺序和边界，
目标损失只作用于约定的当前译句或整篇目标，不能无意重复计算同一监督。
若加入历史译文，训练读取人工译文而运行读取机器译文的差异也需说明；
实时场景不能使用未来句子，离线场景则可以在任务允许时双向读取。

另用教学篇章“海棠研究所发布了新通知。该所还说明了报名规则。”
预先确认机构译名为Haitang Institute。
句级系统只读第二句时，可写出“The institute also explained the registration rules”，
但不知道应同前句哪一机构对齐；读取前句及实体表后，
可统一为“Haitang Institute issued a new notice.
The institute also explained the registration rules.”
若检索上下文含同名但不同地点的机构，仍须以身份而非名称字符串确定对应。

省略也会影响目标语言的必填成分。材料“林岚提交了申请。随后收到了确认。”
第二句省略的接收者由前句恢复为林岚，而不是发出确认的机构。
单看第二句生成“An acknowledgement was received”保留了被动事件，
却没有恢复人物；若人物身份已确认，上下文版本可写
“Lin Lan then received an acknowledgement”。
对这类现象需要逐项对比译文，不能用句级平均分断言所有指代已经正确。

固定两句上下文时，句级路线分别生成后可用实体表协调名称，
但若前一步译错主客体，只替换名称不足以修复关系；
分段协调还须回查对应源句。整篇生成能同时读取全部句子，输出后却仍须按段对齐，
防止漏句或把信息移到错误段落。保存“源实体、目标名称、代词候选、出处”的记录，
可以在受约束解码和后核验之间共享；词汇约束见第~\ref{sec:nlp-generation-constraints}节。

\subsubsection{DocRepair目标语篇章修订}
\label{sec:nlp-generation-docrepair}

有时只能取得已经完成的句级译文，或希望利用丰富的目标语单语篇章训练连贯性修订。
DocRepair不直接编码源语上下文，而把有缺陷的目标语篇章改成连贯版本。
训练先将真实目标语文档逐句译到另一语言，再逐句译回，
以往返翻译后的句子集合为输入，原目标语篇章为监督输出
\scite{voita-etal-2019-context}

逐句往返会丢失部分跨句线索，给名称、指代和衔接修订制造训练样本；
修订模型用条件最大似然学习恢复原篇章。运行时读取句级翻译器的实际输出，
不再把训练用的原目标语文档当作可访问参考。
往返噪声与实际翻译错误的分布可能不同，仍需用独立真实译文检查收益和过改。

教学输入为“Haitang Institute issued a new notice.
Begonia Institute said he would update it.”
在已确认两句指同一机构、第二句源意为“该所表示将更新通知”的条件下，
修订候选为“Haitang Institute issued a new notice.
The institute said it would update the notice.”
名称、机构代词和it的指代都更一致。与上一小节不同，
DocRepair本身只从目标语上下文预测这些修订；
本例给出的源语确认属于独立核验信息，不能说模型已经据此判断。

如果第二句原文实际指另一机构，强行统一名称就会过改；
若it可能指通知或报名表，只读目标语也未必能恢复正确对象。
目标语篇章修订因此不能替代源语证据。术语一致、篇章适切与事实正确应分别核验，
尤其不能把“修订后读起来更连贯”直接记为翻译准确性提高。

\section{生成与转换结果的综合比较}
\label{sec:nlp-generation-comparison}

前面几类任务都能调用条件生成器，却不能共用一份不加区分的验收答案。
通知译文要保留校内外参加者的差别，校外指南摘要可以不写校内规则，
纠错则不应借机删去任何一类参加者。表~\ref{tab:nlp-generation-contracts}把这些差别
落实到给定内容、允许变化和核验依据。表中仍使用本章通知；开放写作另外给出虚构许可，
纠错输入则是人为加入错字的通知。长度、模型规模和候选预算固定后，才能比较代表实现
是否满足同一个任务约定，不能把不同用途的输出仅按流畅程度排成一列。

\begin{table*}[htbp]
  \centering
  \small
  \setlength{\tabcolsep}{3pt}
  \renewcommand{\tabularxcolumn}[1]{>{\raggedright\arraybackslash}p{#1}}
  \begin{tabularx}{\linewidth}{@{}>{\raggedright\arraybackslash}p{13mm}XXXXX@{}}
    \toprule
    任务 & 给定内容 & 必须保留 & 允许改变或新增 & 代表实现 & 核验依据\\
    \midrule
    翻译 & 通知、语向、名称表 & 实体、时间、角色与条件 & 目标语措辞、语序和句界 &
    条件翻译、术语约束 & 源句对齐、术语与原意\\
    \addlinespace
    摘要 & 通知、校外读者、两句预算 & 当前安排、校外截止条件 &
    删去与用途无关的校内规则，合并表达 & 查询选择、MMR、规划生成 &
    必选内容覆盖、逐条来源支持\\
    \addlinespace
    改写 & 通知、正式文体要求 & 原陈述及条件范围 &
    措辞和句法，不增加承诺 & 条件改写、局部编辑 &
    改写前后含义与目标文体\\
    \addlinespace
    纠错 & 含“请与周五到场”的通知 & 正确信息、作者意图 &
    将误用的“与”改为“于” & 混淆集、编辑标签 &
    参考编辑、遗漏与过改\\
    \addlinespace
    简化 & 通知、目标读者、可略背景 & 日期、否定及参加条件 &
    易读词、短句、明确指代 & 规则拆句、属性控制 &
    信息保持、操作评价与读者理解\\
    \addlinespace
    \shortstack[l]{数据到\\文本} & 五条当前安排记录 & 用途要求的记录值与绑定 &
    表述顺序和合句，不补未知事实 & 槽位模板、内容规划 &
    记录覆盖、重复与无依据新增\\
    \addlinespace
    开放写作 & 题目、人物、已知安排及虚构许可 & 必须真实的安排、人物一致性 &
    许可范围内的情节与描写 & 先规划后写作、受控续写 &
    设定一致、用途、创作质量\\
    \bottomrule
  \end{tabularx}
  \caption{生成任务的信息变化与验收。每行的允许变化由用途确定，不是该任务永久不变的规则；
  例如完整摘要与校外指南的必选内容可以不同。代表实现对应本章已经走通的输入和恢复流程，
  不能仅凭方法名称判断是否合格。}
  \label{tab:nlp-generation-contracts}
\end{table*}
\FloatBarrier

\subsection{信息保持与忠实性}
\label{sec:nlp-generation-faithfulness}

要检查“忠实”，需要先确定忠实于什么。翻译以源文原意为依据，摘要以选定来源及用途为
依据，编辑还要检查修改的必要性。开放故事中的虚构人物可以说虚构的话，
但引用的真实活动安排仍不能改成另一天。
把实体、日期和数量规范化后逐项比对，是一个可解释的规则基线：
“周五”在本章背景中对应2025年4月18日，“4月16日18:00前”不能只剩一个日期。
规则还须保留字段角色，否则把报名时间与开会时间互换仍可能通过数值集合检查。

来源并未说出的信息也不应自动当作虚假。通知没有人数，生成“预计50人到场”缺少支持；
生成“活动仍于4月17日举行”则与改期通知相矛盾。
两者都不能作为已核实事实交付，但定位原因和修订方法不同。
下面两条学习式核验路线分别把问题化为句对判断和问答对照，人工核验用于独立检查它们
有没有漏掉关键错误，而不只是复核高分候选。

\subsubsection{SummaC的句对推断核验}
\label{sec:nlp-generation-summac}

把整篇来源和摘要一次性交给句对模型，可能超过模型训练时见到的粒度和长度。
SummaC先分句，寻找每个摘要句在来源中的支持，再汇总为文档分数。
设来源句为$x_1,\ldots,x_M$，摘要句为$y_1,\ldots,y_N$，
句对模型给出蕴含概率$E_{ij}=p(\text{蕴含}\mid x_i,y_j)$，
并可另外保存矛盾与中立分数。蕴含、矛盾和中立的分类基础见前章的句对推断。
SummaC-ZS对每个摘要句取最佳支持，再对摘要句平均：
\begin{equation}
  s_{\mathrm{ZS}}=\frac{1}{N}\sum_{j=1}^{N}\max_{1\leq i\leq M}E_{ij}.
  \label{eq:nlp-generation-summac}
\end{equation}
这个聚合本身没有额外可学习参数；SummaC-Conv则把每列分数离散成直方图，
用可学习的卷积模块获得该列的一致性分数，再平均。两者不能混称为同一个训练过程
\scite{laban-etal-2022-summac}

把通知第一句作为$x_1$，参加规则作为$x_2$，候选两句为
$y_1$：“读书会改至4月18日举行”，$y_2$：“预计50人到场”。
假设教学蕴含矩阵为
\[
  E=\begin{pmatrix}0.95&0.05\\0.10&0.02\end{pmatrix}.
\]
按列取最大值得$(0.95,0.05)$，平均为$0.50$，低分定位到人数句。
若用四个等宽区间构造直方图，第一列在最低、最高区间各有一个元素，
第二列两个元素都落在最低区间；这就是Conv可以读取的分布信息，
但没有训练好的卷积参数时，不能继续编造一个“Conv实测分数”。
假定独立人工标注开发集选出的阈值为$0.80$，该候选会进入待核查集合；
阈值数值仅用于演示，不能不经校准迁移到新语种或长文。

人数句的低蕴含分数不等于高矛盾分数。它对$x_1$完全可能是中立$0.92$、
矛盾$0.03$；若改成“活动仍在4月17日举行”，才可能得到高矛盾分数。
这些仍是教学预测，须回看具体证据，不能让模型概率代替来源本身。
若来源分成“读书会已经改期”和“新日期为4月18日”两句，
单独一句都可能无法充分支持摘要中的完整陈述。可以用相邻句窗口重新核查，
并标记“需要跨句证据”，而不是把分句造成的缺口直接判为幻觉。
逐句最大值还可能忽略来源内部冲突，平均分则可能掩盖唯一但严重的日期错误，
因此应同时保存问题句、对应来源和最低支持项。

\subsubsection{QAFactEval的问答核验}
\label{sec:nlp-generation-qafacteval}

另一种检查方式是从候选中的陈述反向提问，看来源能否给出相同答案。
QAFactEval从摘要选择实体、名词短语等答案片段，围绕片段生成问题；
先检查这些问题在摘要中是否可回答，再到来源寻找答案，比较答案的一致程度。
匹配器可以使用学习式答案评价，而不局限于字符串全等
\scite{fabbri-etal-2022-qafacteval}

沿用含错误人数的摘要，另加正确的第三句
“校外来宾须在4月16日18:00前提交报名信息”。
选择的三个答案分别为“4月18日”“50人”和第三句的完整报名条件。
问题为“读书会现在哪天举行”“预计多少人到场”
和“校外来宾参加前需要在何时完成什么”。
在摘要上回问，可分别恢复三个预定答案，这一步排除了与选中答案无关的坏问题；
到原通知阅读，第一题恢复4月18日，第二题不可回答，
第三题恢复“在4月16日18:00前提交报名信息”。
日期格式统一后，第一、三题匹配，第二题记零。若为教学采用二值匹配，
均值为$(1+0+1)/3=2/3$；这不是QAFactEval真实匹配模型的输出尺度。

更一般地，过滤后保留问题集$\mathcal{Q}$，令$m_q$表示来源答案与候选答案的匹配分，
来源不可回答时将$m_q$置零，再计算
$|\mathcal{Q}|^{-1}\sum_{q\in\mathcal{Q}}m_q$。
若没有问题通过过滤，应报告“无法核验”，不能把空集合平均值当作满分。
QAFactEval中使用摘要可回答性过滤和学习式匹配的具体组合见原论文；
问答系统如何训练与恢复答案由第~\ref{chap:nlp-question-answering-dialogue}章展开。

这里至少有三种可能的错误来源。人数确实是候选无依据新增；
若问题生成成“50人何时到场”，就混入了未经证实的前提；
若来源阅读器把4月16日当成活动日期，失败发生在阅读步骤。
核验日志因此要保留答案片段、问题、摘要回答、来源回答和出处，才能区分原因。
问题选择也可能完全漏掉14:00或海棠楼201，得分高不表示完整性已经通过。
仅问“什么时候”而没有问“谁需要报名”，还会漏掉把校外规则强加给校内师生的错误。
重要信息的提问覆盖率与答案匹配应分开统计；更细的原子事实支持核验见下一章，
其文献入口FActScore列于本章文献说明。

\subsection{写作条件与资源取舍}
\label{sec:nlp-generation-constraints}

“两句”“正式”“必须使用术语表”“五秒内返回”影响不同环节。
提示给生成器条件，编辑修改已有文本，筛选只接受合格候选，解码限制则在候选形成时排除
非法路径。若规定必须保留完整报名条件，同时限制为十个汉字，可能根本没有可行答案；
此时应报告约束冲突，不能默默放弃事实条件。
下面固定材料和最大生成步数，分别观察约束满足、内容正确和搜索开销。
运行预算还应计入属性模型、解析器、修订次数与人工等待，而不只统计生成器调用次数。

\subsubsection{Grid Beam Search词汇约束}
\label{sec:nlp-generation-gbs}

只在束搜索结束后筛选必含术语，可能发现整个候选集合都漏了一个低概率名称。
\term{网格束搜索}（Grid Beam Search, GBS）在生成步数之外增加约束进度维度，
给不同进度分别保留候选，使“已经完成部分术语”的前缀不必立刻与所有普通前缀竞争
\scite{hokamp-liu-2017-lexically}

一个状态应保存文本前缀、累计分数、已完成约束集合，以及当前短语内部的位置。
处于开放状态时，可以生成普通词，也可以启动一条尚未完成的约束；
进入短语后，只能续接该短语的后继词元，直到短语完成再回到开放状态。
网格坐标$(t,j)$表示已生成$t$个词元、已经推进$j$个约束词元，
每个格保留至多$b$项。$j$不是充分状态：同样推进两个词元，
可能刚完成第一条短语，也可能仍在另一条更长短语内部。
只有所有必含短语均已完成、又生成EOS的候选才能交付。

用两个双词元术语演示：$ab$代表Haitang Hall，$cd$代表external guests，
$o$代表其他文字片段。它们只是分词后的符号，本例检查术语状态，不把符号串当成完整译文。
首步概率为$p(o)=0.50,p(a)=0.30,p(c)=0.15,p(\mathrm{EOS})=0.05$。
普通前缀之后的分布为$(o,a,c,\mathrm{EOS})=(0.40,0.25,0.25,0.10)$；
以$a$结尾时$(b,o,\mathrm{EOS})=(0.60,0.30,0.10)$；
以$c$结尾时$(d,o,\mathrm{EOS})=(0.80,0.10,0.10)$；
以$ab$结尾时$(c,o,\mathrm{EOS})=(0.50,0.30,0.20)$；
以$cd$结尾时$(\mathrm{EOS},o)=(0.90,0.10)$，未列词元概率为零。
按状态选用对应行，不能在词汇屏蔽后改写原模型累计概率。

取每格$b=1$，平局按$a<c<o$，最多五步且EOS占一步。
第一步，$j=0$保留$o:0.50$；$j=1$在$a:0.30,c:0.15$中保留$a$，
其内部状态为“等待$b$”。第二步得到
\[
  \begin{array}{c|ccc}
    j&0&1&2\\ \hline
    \text{保留项}&oo:0.20&oa:0.125&ab:0.18
  \end{array}
\]
$oa$和$oc$在$j=1$同分；$ab$完成第一条短语，但尚不能接受EOS。
后续一条可行轨迹为
$ab\rightarrow abc\rightarrow abcd\rightarrow abcd\,\mathrm{EOS}$，
累计概率依次为$0.18,0.09,0.072,0.0648$。
在$abc$处，第二条短语尚待$d$，计数为3而非“两条约束已完成”。
本例保留的其他路径来不及在第五步前补齐术语及EOS。

两条不重叠双词元短语至少需要四个正文词元和一个EOS。
把预算改成四步，任何路径都不可能完成；若再要求“必须含$ab$”且“禁止$b$”，
则即使增加预算也无解。普通束宽一会沿$o,oo,\ldots$走下去；
生成后筛选只能丢弃它，无法补回被剪掉的$a$。
GBS在总约束词元数$C$下，每步至多保存$(C+1)b$项，
比相同每格宽度的普通束需要维护更多前缀，但仍不保证找到全局最优可行文本。
即使输出出现external guests，也可能把它错误地放在“无需报名”之前，
所以词汇满足必须与含义核验分开。

\subsubsection{Dynamic Beam Allocation约束搜索}
\label{sec:nlp-generation-dba}

术语一多，GBS为每个进度格分别保留束的开销随约束数增加。
\term{动态束分配}（Dynamic Beam Allocation, DBA）固定总束宽，
在约束进度组之间重新分配名额。原方法的候选并集包含全局高分扩展、
各前缀对未完成约束的扩展，以及各前缀自身的最佳后继；
候选按约束进度分组后，再按分配规则进入总束
\scite{post-vilar-2018-fast}

沿用上一小节的分布与术语，固定总宽度为2。
为了明确演示“哪一组得到名额”，本例采用如下教学分配：
有两组时各一项；有三组而只有两个名额时，给最高、最低进度组各一项，
组内按模型分数选优，空名额转给其他非空组。这一约定说明固定容量的取舍，
不声称复现原论文所有配额调整细节。
第一步的候选按进度分为$\{o:0.50\}$与$\{a:0.30,c:0.15\}$，
保留$o,a$。第二步的高分与约束扩展并集包含
$oo:0.20,oa:0.125,oc:0.125,ab:0.18$，分到进度0、1、2三组。
两个名额分给$oo$与$ab$，中间进度组这一步没有候选；
GBS则会另外保留$oa$。

两步之后两种方法都还保留成功轨迹的前缀$ab$，
所以在本例中仍有机会完成$abcd\,\mathrm{EOS}$，
但DBA已经失去一条GBS尚保留的备用路径。若未来分布变化，这个差别可能改变最终结果；
固定束宽不是对任意约束数的最优性保证。
每条候选仍需保存短语身份与内部进度。允许普通扩展打断未完成短语时，
DBA必须撤销这段未完成进度，例如$ao$不能继续声称完成了短语$ab$的一半；
以后若再次出现$a$，应重新开始匹配。
已完整出现的短语则继续保留完成标记。

总活动束固定，主要模型前缀数不再按每格宽度乘约束数增长；
构造约束扩展、更新匹配状态及分组仍有约束相关开销。
因此“固定束宽”不等于“所有计算与约束数无关”。
比较时应同时记录总前缀数、候选扩展数、失败率和实际耗时，
并按同一EOS、长度与内容核验规则交付结果。

\subsubsection{FUDGE软属性控制}
\label{sec:nlp-generation-fudge}

正式、礼貌等属性难以写成一个必含词表。FUDGE让一个额外分类器预测
“当前前缀继续写完后，完整文本会不会满足属性”，再用这个预测调整下一词元分布。
训练样本是带完整文本属性标签的序列，把同一个完整文本标签附给其各个前缀，
用二元交叉熵训练前缀分类器；这不是判断前缀本身是否已经写得正式。
生成器可以保持不变，另训属性模型
\scite{yang-klein-2021-fudge}

记要求的属性为$A$，属性模型为$q_{\bm{\phi}}(A\mid h)$。
理想条件概率分解给出如下逐步重加权形式：
\begin{equation}
  \widetilde p(v\mid h,A)=
  \frac{p_{\bm{\theta}}(v\mid h)\,q_{\bm{\phi}}(A\mid hv)}
  {\sum_w p_{\bm{\theta}}(w\mid h)\,q_{\bm{\phi}}(A\mid hw)}.
  \label{eq:nlp-generation-fudge}
\end{equation}
输入材料也属于条件，这里为简洁略去。近似分类器不一定与生成器的真实未来分布校准一致，
因而公式实现的是一种控制近似，不是属性满足的证明。

给定开头“校外来宾”，三个教学候选片段为“须提交”“记得交”“无需提交”，
基础概率为$(0.40,0.50,0.10)$，预测最终正式的概率为$(0.90,0.20,0.95)$。
乘积为$(0.36,0.10,0.095)$，总和$0.555$，归一化为
$(0.64865,0.18018,0.17117)$。束宽一时选择从口语的“记得交”
变成“须提交”；若使用采样，则仍须固定温度、随机数和总生成预算。
“无需提交”同样很正式，控制后它的相对概率甚至上升，
却与通知要求矛盾。属性得分因此不能代替内容检查。

多个属性模型可以尝试相乘，但只有在相应条件独立假设下才能作相应概率解释；
“极短”与“完整保留长条件”还可能直接冲突。
实际实现若只对基础分布的前若干词元计算属性分数，也会改变可搜索集合。
比较FUDGE与硬术语约束时，应区分软属性准确率、硬约束满足率和事实保持，
并计入每一步额外的属性预测开销。

\subsubsection{有限状态与语法约束解码}
\label{sec:nlp-generation-grammar}

若交付的是带字段的文本，最直接的要求可能是“能被解析器读入”。
可以用有限状态机规定固定字段次序、日期字符模式及分隔符，
或者用带栈的语法状态处理嵌套结构。每生成一个模型词元，就尝试把它解码成字符或字节串，
交给当前规则状态推进；无法成为任何合法后续前缀的词元被屏蔽。
一个模型词元可能跨过引号、字段名和冒号等多个终端符，不能简单地把模型词元当成语法终端。

本例只允许两个固定字段，目标形如
\begin{center}
  \code{\{"date":"2025-04-18","room":"201"\}}
\end{center}
为便于追踪，把它拆成五个教学片段：
$q_0$接受\code{\{"date":"}后进入日期状态$q_1$；
$q_1$接受形如\code{2025-04-DD}的日期字符后进入$q_2$；
$q_2$接受\code{","room":"}后进入房间状态$q_3$；
$q_3$接受三位数字后进入$q_4$；
$q_4$接受\code{"\}}后进入完整状态$q_5$，只有$q_5$才允许EOS。
真实词表中的片段划分不同，只要逐字符状态推进相同即可。
日期模式还可以检查月份和日范围，但日历合法仍不等于通知日期正确。

在$q_2$，缺少引号的\code{,room:}被屏蔽，正确分隔片段保留；
若剩余候选都非法，或预算已不足以闭合结构，应返回解码失败。
在$q_1$，\code{2025-04-17}和\code{2025-04-18}都符合格式，
前者却与当前活动日期不符；把它恢复成完整JSON，也只能通过语法检查。
输出取4月18日、房间201并正确闭合后，还须检查活动编号、来源版本以及这两个值的绑定。
解析成功、值正确和内容覆盖是三项不同的验收。

有限状态机适合固定模式，任意深度的嵌套通常需要语法栈或另设深度上限。
规则由接口模式或人工确认的约束提供，不从测试参考临时提取；
状态更新和合法词元集合的构造也占用时间与内存。
若硬词汇、格式和长度约束联合后不可满足，应尽早暴露冲突，
不能以返回一个可解析但事实错误的对象作为降级成功。

\subsection{自动指标、人工评价与质量估计}
\label{sec:nlp-generation-evaluation}

生成任务往往有多种正确写法。参考文本能展示其中一些，却未必覆盖所有合法表达；
来源则规定哪些内容有依据，却不直接给出唯一的好句子。
评价需要同时说明检查对象和依据：参考重合检查表达接近程度，编辑指标检查修改，
学习式指标预测人工质量，人工评审再按用途定位错误。
流畅度、偏好、事实支持与完整性应分别报告，不能把它们都改名为“准确率”。
本节的小例只演示计算，不用不同指标的数值大小直接比较质量。

\subsubsection{BLEU、ROUGE与chrF参考重合}
\label{sec:nlp-generation-overlap}

BLEU用词$n$元组的重合衡量候选与参考的接近程度，并惩罚过短输出。
为防止重复一个常见词刷高分，它对候选计数作截断：
每种$n$元组最多匹配参考中的出现次数，多参考时取各参考计数的最大值。
语料级精确率$p_n$先在句内截断，再跨句累加匹配数和候选总数。
设候选总长度为$c$，按固定规则选出的有效参考总长度为$r$，则
\begin{equation}
  \mathrm{BLEU}_N=\mathrm{BP}\exp\!\left(\sum_{n=1}^{N}w_n\log p_n\right),
  \quad
  \mathrm{BP}=
  \begin{cases}
    1,&c\geq r,\\
    \exp(1-r/c),&0<c<r.
  \end{cases}
  \label{eq:nlp-generation-bleu}
\end{equation}
其中$w_n$非负且和为一；空候选另记零。
原始BLEU（\citeyear{papineni-etal-2002-bleu}）的完整出处见本章文献说明。

取单参考“the meeting is on april 18”，候选“the meeting is on april 17”，
统一小写、按空格分词、不含标点。六个一元组匹配五个，五个二元组匹配四个，
所以$p_1=5/6,p_2=4/5$。采用$N=2,w_1=w_2=1/2$、不平滑的教学口径，
长度相等使$\mathrm{BP}=1$，$\mathrm{BLEU}_2=\sqrt{2/3}\approx0.8165$。
日期错误只损失少量重合，得分仍高。若候选只有“the meeting”，
两阶精确率都是1，长度惩罚却为$\exp(1-6/2)=e^{-2}\approx0.1353$。
若候选为“the the meeting”，其中两个the只有一个能匹配，
一元组精确率为$2/3$，不是1。

标准系统比较常采用四阶BLEU；短句中某阶没有匹配会使不平滑的几何平均为零，
需要明确是否使用平滑和有效阶数。语料BLEU也不等于逐句BLEU的平均。
分词器、大小写、参考数、平滑和汇总方式改变后，分数不能直接横比。
合法改写“the event takes place on april 18”保留事件和日期，
却改变多个词组；仅有一条参考会低估这类表达，而增加参考仍不能保证穷尽合法译法。

ROUGE是一组参考重合指标，在摘要中常用于检查参考信息是否被覆盖。
ROUGE-L用\term{最长公共子序列}（longest common subsequence, LCS），
允许匹配词之间隔着其他词，但保持顺序。对单个候选和参考，设LCS长度为$L$，
则$P=L/c,R=L/r$，并计算
\[
  F_\beta=\frac{(1+\beta^2)PR}{\beta^2P+R}.
\]
$\beta$调节召回的权重。本节明示使用$\beta=1$，而不是把历史上的召回口径或
所有ROUGE实现都当作F1；原始定义见ROUGE论文
\scite{lin-2004-rouge}

上例错误日期候选与参考的LCS是“the meeting is on april”，长度5，
故$P=R=F_1=5/6\approx0.8333$。它同样没有把日期错误视为一票否决。
多句摘要还需说明采用逐句、整篇还是跨句汇总的变体，以及句界处理；
只写“ROUGE”不足以重现评价。

词边界不稳定或形态变化丰富时，可以改看字符重合。
chrF对各阶字符$n$元组分别计算精确率$P_n$和召回率$R_n$，
先取$\bar P=N^{-1}\sum_nP_n,\bar R=N^{-1}\sum_nR_n$，
再代入上面的$F_\beta$，而不是先对各阶F值求平均
\scite{popovic-2015-chrf}

取参考“4月18日”、候选“4月17日”，不计空格、按Unicode字符切分。
五个一元字符中匹配四个，四个二元字符中只有“4月”“月1”匹配，
所以$P_1=R_1=4/5$，$P_2=R_2=1/2$。
教学取$N=2,\beta=2$，得到$\bar P=\bar R=0.65$、chrF$=0.65$。
这是明确的二阶例子，不是工具包默认多阶配置的标准报告值。
字符匹配减轻了某些分词问题，却不会理解数字、否定或角色的业务含义；
系统比较仍须固定阶数、空格、大小写、Unicode规范化和聚合口径。

\subsubsection{SARI、M2与ERRANT编辑评价}
\label{sec:nlp-generation-edit-metrics}

简化不能仅奖励与参考相似，还要知道相对于来源做了什么。
SARI分别评价增加、保留和删除的$n$元组：增加与保留用F1，删除用精确率，
避免只因参考允许删掉很多词就奖励不加选择的删除
\scite{xu-etal-2016-optimizing}

在单参考、集合计数的教学口径下，记来源、候选、参考的$n$元组集合为$X_n,Y_n,R_n$。
增加操作比较$Y_n\setminus X_n$与$R_n\setminus X_n$，
保留操作比较$Y_n\cap X_n$与$R_n\cap X_n$，
删除操作比较$X_n\setminus Y_n$与$X_n\setminus R_n$。
对前两类，精确率的分母是候选操作数，召回率的分母是参考操作数；
删除只取正确删除数除以候选删除数。
每阶三项平均，再对一至四阶平均：
\begin{equation}
  \mathrm{SARI}=\frac{1}{4}\sum_{n=1}^{4}
  \frac{F_{\mathrm{add},n}+F_{\mathrm{keep},n}+P_{\mathrm{del},n}}{3}.
  \label{eq:nlp-generation-sari}
\end{equation}
本例沿用原文零分母记零的约定；多参考时还涉及参考频次权重，不能把参考简单拼接，
也不能仅计算一阶就称为完整SARI。

来源为“the previously announced meeting starts today”，
候选为“the meeting starts today”，参考为“the meeting begins today”。
任务预先允许略去“此前已宣布”的背景，三句按空格分词。
一阶增加中，候选没有新增词，参考增加begins，F值为0；
候选保留the、meeting、starts、today四词，参考要求保留其中三词，
所以$P_{\mathrm{keep},1}=3/4,R_{\mathrm{keep},1}=1$，F值为$6/7$；
候选删除previously与announced，二者都属于参考删除，删除精确率为1。
二阶候选只新增the meeting，参考新增the meeting、meeting begins、begins today，
故增加精确率为1、召回率为$1/3$、F值为$1/2$。
候选保留的meeting starts、starts today都不在参考保留集合中，保留F值为0。
全部阶数的操作数和得分为
\begin{center}
  \small
  \begin{tabular}{@{}crrrccc@{}}
    \toprule
    阶 & 增加数 & 保留数 & 删除数 & 增加F1 & 保留F1 & 删除P\\
    \midrule
    1 & 0/1/0 & 4/3/3 & 2/3/2 & 0 & $6/7$ & 1\\
    2 & 1/3/1 & 2/0/0 & 3/5/3 & $1/2$ & 0 & 1\\
    3 & 1/2/0 & 1/0/0 & 3/4/3 & 0 & 0 & 1\\
    4 & 1/1/0 & 0/0/0 & 3/3/3 & 0 & 0 & 1\\
    \bottomrule
  \end{tabular}
\end{center}
每个操作数单元按“候选操作数／参考操作数／匹配数”排列。
三阶候选新增the meeting starts而参考新增的两项均不相同；
四阶两句的完整短语也不同，所以对应增加项为零。
最终$\mathrm{SARI}=(13/21+1/2+1/3+1/3)/4=25/56\approx0.4464$。
starts与begins在这里可以表达相近意思，参考操作仍偏好后者，说明操作匹配不等于语义判决。
背景是否允许删除也必须由任务决定，不能反过来由SARI高低决定。

纠错更关心具体编辑是否必要。M2从来源到输出的编辑格中保留可能的最短对齐，
加入跨越多个基本操作的短语编辑边，再选择最能匹配参考标注的编辑路径。
同一变化可能是“替换一个短语”，也可能分成多个替换与插入，
因此M2不等同于随意回溯一条Levenshtein最短路径
\scite{dahlmeier-ng-2012-better}

取来源“He go to school yesterday with book .”，零起始词坐标下，
参考编辑为$e_1: [1,2)\mapsto\text{went}$和$e_2:[6,6)\mapsto\text{a}$，
恢复结果为“He went to school yesterday with a book .”。
固定这份编辑标注，候选甲执行两项，匹配2、预测2、参考2，精确率和召回率都是1；
候选乙只把go改为went，匹配1、预测1、参考2，$P=1,R=1/2$；
候选丙在甲上还把He改成The boy，新增了来源未支持的人物描述，
匹配2、预测3、参考2，$P=2/3,R=1$。
若采用偏重精确率的$F_{0.5}$，乙得$5/6$，丙得$5/7$；
若采用F1，乙得$2/3$，丙得$4/5$，排序恰好不同。
原始M2论文报告F1，后续纠错任务中的$F_{0.5}$是另外指定的评价协议。
一项严重错误是否可以用其他正确编辑补偿，也应另按用途检查。

ERRANT把“哪里改了”进一步连接到“改的是什么错误”。
它用词元和词性辅助的Damerau--Levenshtein对齐，再按语言相关规则归并编辑，
依据来源与目标的语言特征分类，而不是给任意编辑距离路径贴一个名称
\scite{bryant-etal-2017-automatic}

在上述英语例中，go到went可归入动词时态替换，缺a对应限定词缺失；
按这两个参考类别计数，乙的时态召回为1、限定词召回为0，
比总分更直接地指出遗漏。He到The boy的额外改写仍须检查依据，
不能因语法合法就算正确修订。实际类别及边界以固定版本的分析器和规则输出为准，
本例没有声称运行了ERRANT工具。
中文需另定分词、编辑边界和错误类别，多参考要明确可匹配的合法修改集合，
不能直接套用英文词形与限定词分类。

\subsubsection{BERTScore语义匹配}
\label{sec:nlp-generation-bertscore}

参考把starts改成begins，字面操作指标就可能扣分。BERTScore改用上下文词元向量，
让候选中的每个词元寻找最相似的参考词元，再从参考向候选反向匹配。
设归一化后的候选向量为$\bm h_i$，参考向量为$\bm r_j$，
无权重版本计算
\begin{equation}
  P_{\mathrm{B}}=\frac{1}{|y|}\sum_i\max_j\bm h_i^\top\bm r_j,\qquad
  R_{\mathrm{B}}=\frac{1}{|r|}\sum_j\max_i\bm h_i^\top\bm r_j.
  \label{eq:nlp-generation-bertscore}
\end{equation}
再由二者计算F1。这是双向最大相似匹配，不是一对一词对齐；
多个词元可以匹配同一个参考词元。表示提取与注意机制见第~\ref{chap:transformer}章，
指标原始定义见章末列出的BERTScore文献。

令参考词元为“甲／帮助／乙”，候选为“甲／协助／乙”。
用四维标准单位向量构造可复核的教学表示：
参考依次取$\bm e_1,\bm e_2,\bm e_3$，候选取
$\bm e_1,0.8\bm e_2+0.6\bm e_4,\bm e_3$。
中间向量的平方范数为$0.8^2+0.6^2=1$，所以得到的确是可实现的余弦矩阵，
对角线为$(1,0.8,1)$，其余为零。双向最大值相同，
$P_{\mathrm B}=R_{\mathrm B}=F_1=2.8/3=14/15\approx0.9333$，
而一元词重合只有$2/3$。它说明向量相似可以给近义替换部分匹配，
不是说某个预训练模型实际产生了这些数值。

再把候选改成“乙／帮助／甲”。若教学表示仍分别靠近
$\bm e_3,\bm e_2,\bm e_1$，矩阵只是交换行，两向最大值全部为1，
但帮助者与被帮助者已经互换。真实上下文模型可能给出不同表示并降低得分，
指标形式却没有保证它一定识别角色错误。数字和否定也应独立核对。
采用逆文档频率权重时，分别把两侧平均改为各自权重归一化的加权平均；
基线重标度还会改变报告尺度。比较必须固定模型、抽取层、分词器、
IDF语料和重标度设置，参考不足与训练偏差也不能由语义向量自动消除。

\subsubsection{COMET与无参考质量估计}
\label{sec:nlp-generation-comet}

词重合和最大向量匹配都没有直接学习“人工认为哪里差、差多少”。
COMET使用人工质量标签训练评价模型，本节取其回归实现说明输入和计算，
不把所有后续版本都视为同一个网络。训练样本为源句$x$、候选$y$、
参考$r$及质量标签$q$；跨语言编码器分别提取句向量
$\bm s,\bm h,\bm r$，再组合候选与来源、参考的逐元素乘积及绝对差
\scite{rei-etal-2020-comet}

原始回归实现把
$[\bm h;\bm r;\bm h\odot\bm s;\bm h\odot\bm r;
|\bm h-\bm s|;|\bm h-\bm r|]$
交给前馈回归器，预测$\widehat q$。
来源向量虽不单独作为一个拼接块，仍通过乘积和差值参与预测。
以均方误差$n^{-1}\sum_i(\widehat q_i-q_i)^2$训练，
标签可来自不同人工协议，但必须固定其定义与尺度。
若教学标签为$(1.0,0.2)$，预测为$(0.8,0.5)$，
这一步损失为$(0.04+0.09)/2=0.065$。
这些数字仅演示回归损失，不是COMET在本章文本上的实测质量。

例如源文补充记录为“资料袋共12份”，正确候选为“There are 12 information packs”，
错误候选把12改成21。两者与同一参考都可能有很高的词重合；
读取来源的评价器有机会根据源译差别降低第二条质量预测，
但预测模型未必可靠识别数字交换，独立数量核验仍应直接判错。
系统级评价可在固定测试集上平均句级预测，同时按文档、语向和现象报告错误；
整体均值不能证明所有关键句合格。

没有参考时，BLEU等参考重合分数没有计算依据。
\term{无参考质量估计}（quality estimation, QE）改用源句与译文预测质量，
仍需在训练时取得人工或规定的质量监督。
不能简单删去已训练参考型模型的参考字段，就称它为同等有效的QE模型。
例如官方接口中的\code{wmt22-comet-da}需要来源、译文和参考，
\code{wmt22-cometkiwi-da}只需来源和译文；
这两个版本的输入、标签与分数尺度需分别记录，官方接口出处见文献说明。
QE可用于把疑似低质译文送人工复核，不能把某个分数解释为“译文正确的概率”。
训练语向之外、陌生领域、术语密集或篇章依赖的输入，应另查校准和错误分布。

\subsubsection{MQM错误标注与人工比较}
\label{sec:nlp-generation-mqm}

总分不能直接告诉译者怎样修订。MQM提供按错误类别和严重程度组织人工标注的框架，
让评价者指出误译、遗漏、术语、流畅性等问题及其影响。
实际研究还要规定类别子集、严重程度、标注单位、来源与上下文、评价者资格、
复评方式以及如何把错误汇总成分数；这些设置不是一个缩写就能确定的。
专家、错误与篇章上下文的评价研究见章末文献说明。

固定源句“校外来宾须在4月16日18:00前提交报名信息”，
候选甲为“Guests are welcome to join”，候选乙为
“External guests must, before 18:00 on April 16, submit registration information”。
甲简短流畅，却没有让校外来宾知道必须完成什么；乙较正式，但保留了角色、动作和截止条件。
只问“更喜欢读哪句”与问“能否据此正确参加活动”，可能得到不同判断。
评价界面应同时提供同一份来源与用途，随机交换候选左右位置，
分别收集偏好、错误位置与用途是否满足，不能拿无来源偏好代替准确性标注。

为演示汇总，预先规定一次关键参加条件遗漏记为一个严重错误，
轻微、严重、致命错误权重分别为1、5、25；这只是本例的项目配置，
不是MQM唯一的标准权重。两名评价者若确认甲有这一处严重遗漏，
甲罚分为5，乙无已识别错误则为0。
若甲某一片段同时引起几个相近错误描述，应按预定规则避免重复计罚；
若用途规定任何关键条件遗漏均不合格，就同时记录甲未通过硬验收，
不能仅靠语料平均罚分把这次失败掩盖。

人工判断也有不确定性。可以报告标注分歧、复评后的错误类型，
并按文档而非把同一文档的相关句子当独立样本，重采样估计系统差异。
一个纯教学的配对文档罚分差为$d=(5,0,5,0)$，定义正值表示甲更差。
若暂假设四份文档独立，平均差为2.5，样本标准差约2.8868，
标准误约1.4434；用自由度3的$t_{0.975}=3.182$得到
区间$2.5\pm3.182\times1.4434\approx[-2.093,7.093]$。
这样的小样本不能据均值断言总体上甲更差；实际错误计数的偏态还可能使正态近似不合适，
应增大样本并报告配对重采样或适合该协议的区间。
这不妨碍对已核实的那一条遗漏作出明确修订。

摘要人工评价还要检查选择与冗余，简化要检查目标读者能否理解，
开放写作则要另外约定情节与创作质量。不能把翻译错误表原封不动当作所有文本的质量定义。
若让模型参与评审，同样需固定材料、提示和标注依据，并用独立人工样本核查偏差；
共同评价原则回用第~\ref{chap:nlp-introduction}章。

\Needspace{8\baselineskip}
\subsection*{文献说明}
\label{sec:nlp-generation-literature}

下面集中给出正文回指的长条目与任务入口。数据集说明只支持相应任务和评价范围，
不表示本章教学例子复现了论文实验。常规方法的出处已在对应段落边注列出。

\begin{fullwidth}
\small
条件摘要所用BART的预训练入口：\fullcite{lewis-etal-2020-bart}
\par\smallskip
Delete--Retrieve--Generate的属性删除、检索与生成设置：
\fullcite{li-etal-2018-delete}
\par\smallskip
JFLEG提供英文流畅性纠错材料：\fullcite{nlp-intro-napoles2017}
\par\smallskip
ASSET提供含多种改写操作的简化参考：\fullcite{nlp-intro-alvamanchego2020}
\par\smallskip
XSum对应高度压缩的新闻摘要任务：\fullcite{nlp-intro-narayan2018}
\par\smallskip
SummEval连接自动分数与多维人工评价：\fullcite{fabbri-etal-2021-summeval}
\par\smallskip
FActScore评估原子事实的来源支持，不能单独证明信息完整：
\fullcite{min-etal-2023-factscore}
\par\smallskip
短语统计翻译的模型与提取入口：\fullcite{nlp-intro-koehn2003}
\par\smallskip
STACL提供前缀可见条件与等待策略：\fullcite{nlp-intro-ma2019stacl}
\par\smallskip
NLLB连接多语数据、模型与语言覆盖评价：\fullcite{nlp-intro-nllb2022}
\par\smallskip
WMT24是固定任务条件下的系统比较入口：\fullcite{nlp-intro-kocmi2024}
\par\smallskip
DiscoX提供中英专业篇章评价材料，不作为实时翻译机制的依据：
\fullcite{nlp-intro-zhao2026discox}
\par\smallskip
BERTScore的词元表示与匹配定义：\fullcite{nlp-generation-zhang2020bertscore}
\par\smallskip
BLEU的截断计数与长度惩罚：\fullcite{papineni-etal-2002-bleu}
\par\smallskip
MQM错误标注、专家评价及上下文影响的研究：
\fullcite{nlp-intro-freitag2021}
\par\smallskip
COMET版本输入以官方文档为准：
\url{https://github.com/Unbabel/COMET}。
这里核对的是参考型与无参考型接口，不据此声称任何版本在本章材料上的实测表现。
\end{fullwidth}

\section{本章小结}
\label{sec:nlp-generation-summary}

通知的译文、日历摘要与易读指南可以使用相同的条件生成器，却承担不同的信息责任。
译文要把原意带到另一种语言，摘要要按用途选择并保留有支持的内容，
编辑则要在允许范围内修改而不引入非预期变化。
因此，任务约定必须在生成前写清材料来源、读者用途、必须保留的信息和允许新增的内容，
不能等结果出来后再用一个分数替它决定标准。

条件最大似然提供学习目标，搜索和采样把概率分布变成候选，
模板、复制、局部编辑与内容规划则给出不同的信息组织方式。
束宽、温度、约束状态和上下文范围改变可探索的候选及成本，
没有任何一项单独保证忠实。输出仍须恢复成可交付的文字或结构，
逐项检查实体绑定、否定、日期、数量和来源支持；无解或预算不足时要明确返回失败。

参考重合、编辑匹配、学习式评价和人工错误标注各提供一部分证据。
它们的作用是帮助比较和定位问题，不是把流畅输出自动确认为正确。
下一章转向问答与对话：材料不一定一开始就已齐备，系统需要围绕问题取得证据，
并在交流中维护上下文、澄清缺失条件。本章已经建立的生成、核验与修订接口
会继续使用，而答案依据和交互状态需要另外建模。
